
In this section we establish universal multiplicative properties for the Grothendieck-Witt space and spectrum, the $\L$-space and spectrum, and the Karoubi-localising variants $\KGW$ and $\KL$ constructed in \S\reffour{subsection:localising-invariants}, compatibly with the respective fundamental fibre squares. The main results, which rely to a large extent on Appendix~\reffour{section:appendix-localisations}, are collected in \S\reffour{subsection:multiplicative-GW}. We then dedicate \S\reffour{subsection:examples} to a comprehensive review of examples, mostly from the world of commutative rings and ring spectra. In \S\reffour{subsection:PeriodicityviaMultiplication} we discuss the compatibility of the above multiplicative structures with Ranciki periodicity and Karoubi's fundamental theorem, while in \S\reffour{subsection:mult-shaneson} we discuss their compatibility with the circle decomposition theorem of \S\reffour{subsection:shaneson}. In \S\reffour{subsection:normal-L} we discuss the induced multiplicative structure on normal L-theory. Finally, in \S\reffour{subsection:automorphism-GW-L} we use the above multiplicative structures to describe the ring of homotopy classes of self-equivalences of $\GW$ and $\L$.



 



\subsection{Day convolution of additive and localising functors}
\label{subsection:multiplicative-GW}%
\label{subsection:multiplicative-L}%
\label{subsection:mult-tate}%


Our construction of the desired multiplicative structures starts by interpreting a lax symmetric monoidal structure on a functor $F\colon \C\to \D$ between symmetric monoidal categories $\C^\otimes$ and $\D^\otimes$ as a commutative algebra structure on $F$ in the Day convolution operad $\Fun(\C, \D)^\Day$, whose underlying $\infty$-category is $\Fun(\C, \D)$ and whose $n$-ary morphisms $(F_1, \dots, F_n)\to G$ are given by the multi-natural transformations $\bigotimes_{i=1}^n F_i(c_i) \to G(\bigotimes_{i=1}^n c_i)$. In the case at hand we have $\C=\Catp$ equipped with its symmetric monoidal structure from~\S \refone{subsection:tensor-product} and \S \refone{subsection:monoidal-structure}, and $\D = \Sps, \Sp$ or $\Grp_\Einf$. Upon restriction to finitary (or more generally accessible) functors, the resulting Day convolution operad turns out to be a symmetric monoidal category. In the presence of such a symmetric monoidal structure, the arguments that we are about to give are well-established, and build on the notion of multiplicative Bousfield localisations. However, the restriction to accessible functors turns out to be unnecessary, because the arguments can be completely performed on an operadic level. We have recorded in Appendix \reffour{section:appendix-localisations} %
the requisite extension of the theory of multiplicative and smashing localisations. In what follows, we shall freely use constructions and results from that appendix.


Taking the target $\D=\Spa$ for illustration, our argument is based on the following key observations:
\begin{enumerate}
\item
The Day convolution $\infty$-operad $\Fun(\Catp,\Spa)^{\Day}$, though not a symmetric monoidal $\infty$-category, still admits a unit object, given explicitly by $\SS[\Poinc]$.
\item
Each of the functors $\GW,\KGW,\L,\KL$ is initial under $\SS[\Poinc]$ with respect to a certain property: being additive, Karoubi-localising, or a combination of one of these with bordism invariance. Each of these properties can be phrased as sending a certain collection of cone diagrams to limit diagrams.
\item
When the collection of cone diagrams is suitably compatible with the tensor product on $\Catp$, such universal approximations of the unit in $\Fun(\Catp,\Spa)^{\Day}$ become themselves the units of the full suboperad spanned by the functors which send the relevant collection of cone diagrams to limit diagrams. In particular, these universal approximations refine to initial commutative algebra objects, hence initial lax symmetric monoidal functors of the given type. %
\end{enumerate}
Replacing the target with $\Grp_{\Einf}$ or $\Sps$ one can also consider in this manner the functors $\GWspace$ and $\Lspace$. On the other hand, one can also replace $\Catp$ with $\Catx$ and consider functors such as $\K,\KK$ or $\Kspace$, thus recovering classical multiplicative properties of $\K$-theory.

To perform our arguments simultaneously for all our examples, we start with symmetric monoidal $\infty$-categories $\C$ and $\D$, %
with $\D$ cocomplete and such that its tensor product preserves colimits in each argument. 
In this situation, the Day convolution operad admits a unit object $\one_\Day = [c \mapsto (\one_\D)_{\Hom_\C(\one_\C, c)}]$, which consequently refines to form the initial lax symmetric monoidal functor $\C \to \D$. 
Let $\cR$ be a collection of cone diagrams $K^{\triangleleft} \rightarrow \C$ (with varying $K \in \Cat$) and denote by $\Fun^\cR(\C,\D) \subseteq \Fun(\C,\D)$ the full subcategory spanned by those functor that take each diagram in $\cR$ to a limit cone. We will refer to such functors as $\cR$-local.
Here are the examples that we have in mind:


\begin{definition}
\label{def:collections-R}%
We define the following collection of cone diagrams in $\Catp$:
\begin{enumerate}
\item We write $\cR_{\add}$ for the collection consisting of $0 \colon \ast \rightarrow \Catp$ and all split Poincaré-Verdier squares.
\item We write $\cR_{\kloc}$ for the collection consisting of $0 \colon \ast \rightarrow \Catp$ and all Poincaré-Karoubi squares.
\item We write $\cR_{\binv}$ for the collection consisting of all cylinder inclusions $\cyl\colon (\C,\QF) \to \Q_1(\C,\QF)$.
\item We write $\cR_{\bord}$ for $\cR_\add \cup \cR_{\binv}$.
\end{enumerate}
\end{definition}


Then, by definition, for an $\infty$-category $\E$, %
a functor $\F\colon \Catp \to \E$ is
\begin{enumerate}
\item $\cR_{\add}$-local if and only if it is additive,
\item $\cR_{\kloc}$-local if and only if it is Karoubi-localising,
\item $\cR_{\binv}$-local if and only if it is bordism invariant, and hence
\item $\cR_\bord$-local if and only if it is additive and bordism invariant, and similarly
\item $(\cR_{\kloc} \cup \cR_{\bord})$-local if and only if it is Karoubi-localising and bordism invariant.
\end{enumerate}

\begin{remark}
\label{remark:binv-is-not-add}%
Let us point out that while in \papertwo we only considered the notion of bordism invariance for additive functors, here we separate additivity and bordism invariance. This leads to a better behaved theory for functors with unstable targets, such as spaces. For example,  
the second part of Corollary~\reffour{corollary:universal-multiplicative-L-space} fails for $\cR_{\bord}$ in place of $\cR_{\binv}$.
\end{remark}



The following result is obtained by applying 
Corollary~\reffour{corollary:multiplicative-day} 
of Appendix \reffour{section:appendix-localisations} 
to the $\infty$-operads $\cO=\Com$ and $\cO=\MCom$. %


\begin{proposition}
\label{proposition:universal-unit}%
Let $\C$ and $\D$ be symmetric monoidal $\infty$-categories %
such that $\D$ is cocomplete and closed, %
and let $\cR$ be a collection of cone diagrams in $\C$ as described above. Assume that $\cR$ is stable under tensor products with arbitrary objects of $\C$, and suppose given a natural transformation $\eta\colon \one_{\Day} \Rightarrow F$ which exhibits $F$ as the initial $\cR$-local functor under $\one_{\Day}$. Then the following holds:
\begin{enumerate}
\item
$F$ carries an essentially unique lax symmetric monoidal structure compatible with $\eta$.
\item
This structure renders $F$ initial in lax symmetric monoidal $\cR$-local functors $\C \to \D$. 
\item
If $\E$ is a cocomplete $\infty$-category equipped with a closed action of $\D$ then any $\cR$-local functor $\C \rightarrow \E$ is a module over $F$ in an essentially unique fashion.
\end{enumerate}
\end{proposition}


It turns out that for some of the functors we consider, notably $\GW$, more is true: not only is $\GW$ a commutative algebra objects with respect to Day convolution, it is actually an idempotent algebra object, and Day convolving with it yields a smashing localisations of $\Fun(\Catp,\Spa)^{\Day}$ into its full suboperad spanned by the additive 
functors. This means in particular that the notion of an additive functor $\G\colon \Catp \to \Spa$ is equivalent to that of $\GW$-module, and every $\G$ admits a universal additive approximation given by Day convolution with $\GW$. The analogous statement also holds for the $\L$-space functor $\Lspace\colon \Catp \to \Sps$ and the $\K$-spectrum and $\K$-space functors on $\Catx$,
for example, but fail for $\KGW,\KK$ and $\KL$, as well as for $\GWspace$ (see Remark~\reffour{remark:not-convolving}). The required sufficient conditions for a smashing localisation to arise are summarised in the following result, %
which is a special case of Proposition~\reffour{proposition:smashing-day} of Appendix \reffour{section:appendix-localisations}. Here, recall that a functor $F\colon \C \to \D$ is said to be \defi{small} if it is left Kan extended from a small subcategory of $\C$. As we establish in Corollary~\reffour{corollary:persist-locally-small}, under the conditions of the following proposition (for instance) small functors are symmetric monoidal objects in $\Fun(\C,\D)^{\Day}$,
that is, Day convolution of a small functor against any functor is well-defined.

\begin{proposition}
\label{proposition:sample-day-smashing}%
Let $\C$ and $\D$ be closed symmetric monoidal $\infty$-categories %
such that $\D$ is cocomplete, %
and let $\cR$ be a collection of cone diagrams in $\C$ as described above. 
Suppose given a natural transformation $\eta\colon \one_{\Day} \Rightarrow F$ 
such that $F$ is small and $\cR$-local,
and such that the following hypotheses hold:
\begin{enumerate}
\item
The collection $\cR$ is stable under tensoring and cotensoring with objects in $\C$. 
\item
The class of $\cR$-local functors is closed under colimits and tensoring with objects of $\D$.
\end{enumerate}
Then $\eta$ exhibits $F$ as the initial $\cR$-local functor under $\one_{\Day}$ if and only if the induced map $G \mapsto F \otimes_{\Day} G$ is an equivalence for every $\cR$-local $G$, in which case
$F$ is idempotent with respect to Day convolution and convolving with it yields a smashing localisation of $\infty$-operads
\[
\Fun(\C,\D)^{\Day} \longrightarrow \Fun^\cR(\C,\D)^{\Day}, \quad G \longmapsto F \otimes_\Day G .
\]
In particular, in this case $\Fun^\cR(\C,\D)^{\Day}$ identifies with the $\infty$-operad of $F$-modules in $\Fun(\C,\D)^{\Day}$.
\end{proposition}


To apply these general results we now verify the relevant hypotheses for the cone diagrams of Definition~\reffour{def:collections-R}. Note that the tensoring and cotensoring mentioned in these hypotheses correspond to the operations $(\C,\QF) \otimes -$ and $\Funx((\C,\QF),-)$ of taking tensor products with, and internal mapping objects from, a given $(\C,\QF) \in \Catp$.  










\begin{proposition}
\label{prop:closure-properties-of-Rs}%
The above defined collections satisfy the following properties:
\begin{enumerate}
\item
\label{item:list1}%
The collections $\cR_{\add}$ and $\cR_{\binv}$ are closed under (co)tensoring with arbitrary $(\C,\QF) \in \Catp$.
\item
\label{item:list2}%
The collection $\cR_{\kloc}$ is closed under tensoring with arbitrary $(\C,\QF) \in \Catp$.
\end{enumerate}
\end{proposition}
\begin{proof}
First we note that $0 \otimes (\C,\QF) = 0 = \Funx((\C,\QF),0)$. 
Combined with Proposition~\reffour{lemma:tensor-preserves-karoubi-squares}, this establishes Part \reffouritem{item:list2}. We now address Part \reffouritem{item:list1}.

To obtain the claim for $\cR_{\add}$ we show that the functors $(\C,\QF)\otimes -$ and $\Funx((\C,\QF),-)$ preserve split Poincar\'e-Verdier squares.
For the latter, since it preserves limits and adjunctions on the level of underlying $\infty$-categories, it suffice to check that it preserves Poincaré-Verdier projections, which follows from Lemma~\reftwo{lemma:and-smashing}.
For the former, we have established in Proposition~\reftwo{lemma:Morita-exact-criterion}\;\reftwoitem{item:preserves-split-PV-sequences} that $(\C,\QF) \otimes -$ preserves split Poincaré-Verdier sequences.
By Proposition~\reftwo{proposition:poincare-verdier-square-char} it hence suffices to show that it also preserves vertically left and right adjointable squares whose vertical legs admit left and right adjoints. This follows again since $\C\otimes -$ preserves adjunctions, so a left or right adjoint to a given vertical leg induces the same kind of adjoint after tensoring with $\C$, with unit and counit induced by the original ones. The resulting Beck-Chevalley transformations are then also induced from the original ones after tensoring with $\C$, and are hence in particular invertible as soon as the original ones were so. 

To see the claim about $\cR_\binv$, we recall that $\Q_1(\D,\QFD)$ is the cotensor $(\D,\QFD)^{\mathcal{P}}$ where $\mathcal{P}$ is the finite poset classifying pushouts as considered earlier.
Under this identification, the Poincaré functor $\cyl$ is the one induced by terminal map $\mathcal{P} \to \ast$.
From the natural equivalence $\Funx((\C,\QF),(\D,\QFD)^{\mathcal{P}}) = \Funx((\C,\QF),(\D,\QFD))^{\mathcal{P}}$ we find that $\cR_{\binv}$ is closed under cotensors. Finally, it follows from Corollary~\refone{corollary:internally-corepresented} that since $\mathcal{P}$ is a finite poset, $(\Spaf,\QF^\uni)^{\mathcal{P}}$ is dualisable in $\Catp$ with dual $(\Spaf,\QF^\uni)_{\mathcal{P}}$. In particular, the natural map
$(\D,\QFD) \otimes (\Spaf,\QF^{\uni})^{\mathcal{P}} \to (\D,\QFD)^{\mathcal{P}} = \Funx((\Spaf,\QF^\uni)_{\mathcal{P}},(\D,\QFD))$ is an equivalence for all $(\D,\Phi)$, implying the remaining claim. 
\end{proof}

\begin{remark}
\label{rem:contensor-does-not-preserve-KS}%
The collection $\cR_{\kloc}$ is not closed under cotensoring with arbitrary $(\C,\QF) \in \Catp$. For instance, $\Funx((\Dperf(\QQ),\QF^\sym),-)$ does not preserve Poincar\'e-Karoubi sequences as it sends the Poincar\'e-Verdier sequence 
\[
(\Dperf(\ZZ)^\mathrm{tors},\QF^\s) \to (\Dperf(\ZZ),\QF^\s) \to (\Dperf(\QQ),\QF^\s)
\]
to the sequence $0 \to 0 \to (\Dperf(\QQ),\QF^\s)$ which is far from a Poincar\'e-Karoubi sequence. 
\end{remark}

Before moving on to our applications of Propositions~\reffour{proposition:universal-unit} and \reffour{proposition:sample-day-smashing},
we record the following:
\begin{remark}
\label{remark:small-functors}%
All of the functors $\Kspace,\GWspace,\Lspace,\K,\GW,\L,\KK,\KGW,\KL$ preserve filtered colimits. Consequently, since $\Catx$ and $\Catp$ are compactly generated (see Proposition~\reffive{proposition:catp-cath-are-compact}), all of these functors are small. Indeed, that $\Kspace$, $\K$ and $\KK$ preserve filtered colimits is well-known. That $\GWspace,\GW,\Lspace$ and $\L$ preserve filtered colimits was noted in Corollaries \reftwo{corollary:filtcolimpresgwspace}, \reftwo{corollary:filtcolimpresgwspectrum}, Theorem~\reftwo{theorem:L-is-localising} and Corollary~\reftwo{corollary:Ladd}, respectively. Finally, for $\KGW$ and $\KL$ this is established in Corollary \reffour{corollary:filtcolimpreskgwspectrum}.
\end{remark}












\begin{corollary}
\label{corollary:GWspace-universal-multiplicative}%
The natural transformation $\Poinc \Rightarrow \GWspace$
refines in an essentially unique manner to a lax symmetric monoidal natural transformation among lax symmetric monoidal functors, making $\GWspace$ into the initial lax symmetric monoidal group-like additive functor $\Catp \to \Sps$.
Furthermore, any group-like additive functor $\Catp \to \E$ valued in a %
cocomplete and finitely complete $\infty$-category $\E$ is a module over $\GWspace$ in an essentially unique manner.
\end{corollary}
\begin{proof}
Recall that group-like additive functors $\Catp \to \Sps$ can be canonically identified with additive functors $\Catp \to \Grp_{\Einf}$ (or, more generally, group-like additive functors $\Catp \to \E$ with additive functors $\Catp \to \Grp_{\Einf}(\E)$). Viewing $\GWspace$ as a $\Grp_{\Einf}$-valued functor, 
we first deduce from Observation~\reftwo{observation:univ-GW-space} that $\GWspace$ is an $\cR_\add$-localisation of the Day unit 
\[
\Fr_{\Einf-\grp}[\Poinc] \in \Fun(\Catp,\Grp_{\Einf})^\Day ,
\]
where $\Fr_{\Einf-\grp}$ denotes the free $\Einf$-group functor, and then apply Proposition~\reffour{proposition:universal-unit}
with $\cR = \R_\add$ using Proposition~\reffour{prop:closure-properties-of-Rs}, obtaining the stated result with $\Fr_{\Einf-\grp}[\Poinc]$ in place of $\Poinc$. We then finish by noting that natural transformations of space-valued functors $\Poinc \Rightarrow \GWspace$ correspond via adjunction to natural transformations of $\Grp_{\Einf}$-valued functors $\Fr_{\Einf-\grp}[\Poinc] \Rightarrow \GWspace$ (and the same holds for lax symmetric monoidal natural transformations; the latter are contractible in both cases since $\Poinc$ and $\Fr_{\Einf-\grp}[\Poinc]$ are the initial lax symmetric monoidal functors to $\Sps$ and $\Grp_{\Einf}$, respectively).
\end{proof}



\begin{corollary}\
\label{corollary:GW-universal-multiplicative}%
\begin{enumerate}
\item
\label{item:GW-universal}%
The natural transformation $\SS[\Poinc] \Rightarrow \GW$ refines in an essentially unique manner to a lax symmetric monoidal natural transformation among lax symmetric monoidal functors, making $\GW$ into the initial 
lax symmetric monoidal additive functor $\Catp \rightarrow \Spa$. Furthermore, any additive functor $\Catp \to \E$ valued in a stable cocomplete $\infty$-category $\E$ is a module over $\GW$ in an essentially unique manner.
\item
\label{item:GW-idempotent}%
The functor $\GW\colon \Catp \to \Spa$ is Day idempotent and Day convolution with it yields a smashing localisation of $\infty$-operads
\[
\Fun(\Catp,\Spa)^{\Day} \locadjL{\GW \otimes_{\Day} (-)} \Funadd(\Catp,\Spa)^{\Day} .
\]
\end{enumerate}
\end{corollary}
\begin{proof}
We deduce from Corollary~\reftwo{corollary:univ-GW-spectrum} that $\GW$ is an $\cR_\add$-localisation of $\SS[\Poinc]$, so the first part follows similarly as above from Proposition~\reffour{proposition:universal-unit}.
To see the second part, note that the collection of spectrum-valued $\cR_\add$-local functors is closed under colimits and tensors with spectra since $\Spa$ is stable. The claim then follows from
Proposition~\reffour{proposition:sample-day-smashing}
since $\GW$ is small by Remark~\reffour{remark:small-functors}.
\end{proof}

\begin{remark}
It turns out that $\GWspace$ is also Day idempotent, both as a space and an $\Einf$-group valued functor, see~\cite{saunier-motives}, but the category it localises onto is not that of additive functors. Instead, a functor is $\G\colon \Catp \to \Grp_{\Einf}$ is invariant under Day convolving with $\GWspace\colon \Catp \to \Grp_{\Einf}$ if and only if for every Poincaré $\infty$-category $\C$ containing an isotropic subcategory $\Lag \hrar \C$ the induced map $\G(\C) \to \G(\Hyp(\Lag)) \times \G(\Hlgy(\Lag))$ is an equivalence, while a functor $\G\colon \Catp \to \Sps$ is invariant under Day convolving with $\GWspace\colon \Catp \to \Sps$ if and only if for every $\C$ and $\Lag \hrar \C$ as before 
the induced map $\G(\C) \to \G(\Hyp(\Lag) \times \Hlgy(\Lag))$ is an equivalence. %
The $\infty$-category of additive functors $\Catp \to \Grp_{\Einf}$ is, in turn, simply not a smashing localisation (otherwise its idempotent would have to be $\GWspace$), and we do not know whether every functor $\Catp \to \Grp_{\Einf}$ admits an additive approximation.
\end{remark}


\begin{corollary}
\label{corollary:kgw-multiplicative}%
The natural transformation $\SS[\Poinc] \Rightarrow \KGW$ refines in an essentially unique manner to a lax symmetric monoidal natural transformation among lax symmetric monoidal functors, making $\KGW$ into the initial 
lax symmetric monoidal Karoubi-localising functor $\Catp \rightarrow \Spa$. Furthermore, any Karoubi-localising functor $\Catp \to \E$ valued in a stable cocomplete $\infty$-category $\E$ is a module over $\GWspace$ in an essentially unique manner.
\end{corollary}
\begin{proof}
By Proposition~\reffour{proposition:universal-KGW} we find that $\KGW$ is a $\cR_\kloc$-localisation of $\SS[\Poinc]$ so that we may apply Proposition~\reffour{proposition:universal-unit}
with $\cR = \cR_\kloc$.
\end{proof}












\begin{remark}\
\label{remark:not-convolving}%
We do not now whether $\KGW$ is idempotent, nor whether every functor $G\colon \Catp \to \Spa$ admits a universal Karoubi-localising approximation $G \to \mathbb{G}$, but we do know that such an approximation cannot be obtained by Day convolving with $\KGW$. Indeed, if that were true then the Day convolution 
\[
\KGW \otimes_\Day \SS[\Hom_{\Catp}((\C,\QF),-)] \simeq \KGW(\Funx((\C,\QF),-))
\]
would be Karoubi-localising, and this fails, e.g., for $(\C,\QF) = (\Dperf(\QQ),\QF^\sym)$ %
by Remark~\reffour{rem:contensor-does-not-preserve-KS}. 
\end{remark}



Before moving on, we will establish the following strengthening of the universal property of the L-space functor established in \papertwo. To avoid confusion, recall from Remark~\reffour{remark:binv-is-not-add} that, unlike in loc.\ cit., here we do not require bordism invariant functors to be additive. %
 
\begin{proposition}
\label{proposition:universal-L-space}%
The functor $\Lspace \colon \Catp \rightarrow \Sps $ is the initial bordism invariant functor under $\Poinc$.
\end{proposition}
\begin{proof}
The functor $\Lspace$ is bordism invariant by Theorem~\reftwo{theorem:L-is-localising}. 
Now suppose that $\G\colon \Catp \to \Sps$ is a bordism invariant functor. From the equivalence 
$\Lspace= \colim_{n \in \Delta\op} y((\Spaf,\QF^\uni)_{\T_n})$, where $\T_n$ denotes the face poset of $\Delta^n$, we deduce that the map
$\Nat(\Lspace,\G) \to \Nat(\Poinc,\G)$
identifies with the projection
\[
\lim_{n \in \Del}\G((\Spaf,\QF^{\uni})_{\T_n}) \to \G(\Spaf,\QF^{\uni}).
\]
To show that this map is an equivalence, by the bordism invariance of $\G$, it suffices to show that each of the Poincaré functors $(\Spaf,\QF^{\uni})_{\T_n} \to (\Spaf,\QF^{\uni})$ is a bordism equivalence. By Corollary~\refone{corollary:internally-corepresented}, this functor is obtained by applying the functor $\Funx(-,(\Spaf,\QF^{\uni}))$ to $(\Spaf,\QF^{\uni}) \to (\Spaf,\QF^{\uni})^{\T_n}$, and it is enough to prove that the latter is a bordism equivalence.
We then observe that this last functor is the inclusion of the homology $\infty$-category $\Hlgy(\Lag_n)$ for an isotropic subcategory $\Lag_n \subseteq (\Spaf)^{\T_n}$ and then use that such an inclusion is a bordism equivalence, see Example~\reftwo{example:isotropic}. To define $\Lag_n$, let $\dM_n \subseteq (\Spaf)^{\T_n}$ be the full subcategory spanned by those diagrams $\vphi\colon  \T_{n} \rightarrow \Spaf$ such that $\vphi(S \cup \{0\}) \rightarrow \vphi(S)$ is an equivalence for every $S$ (equivalently, which are left Kan extended from their restriction to the subposet $\T^0_n \subseteq \T_n$ spanned by those $S$ which contain $0$). We then define $\Lag_n \subseteq \dM_n$ 
to be the full subcategory spanned by those diagrams which additionally satisfy $\varphi(0) \simeq 0$. 
The inclusions $\Lag_n \subseteq \dM_n \subseteq (\Spaf)^{\T_n}$ then both admit right adjoints and one readily checks that for $\varphi \in \dM_n$ there is an equivalence
$
\QF^{\T_n}(\varphi) \simeq \QF(\varphi(0)),
$
so that $\Lag_n$ is indeed isotropic. Furthermore, inspecting the explicit formula for the duality we find that $\Dual^{\T_n}(\dM_n)$ consists of those diagrams whose restriction to $\T^0_n$ is constant, so that
$(\Lag_n)^\perp = \dM_n$ and 
$\Hlgy(\Lag_n) \simeq \dM_n \cap \Dual^{\T_n}(\dM_n)$ contains precisely the constant diagrams.
\end{proof}




\begin{corollary}
\label{corollary:universal-multiplicative-L-space}%
\
\begin{enumerate}
\item
The natural transformation $\Poinc \Rightarrow \Lspace$ refines in an essentially unique manner to a lax symmetric monoidal natural transformation among lax symmetric monoidal functors, making $\Lspace$ into the initial 
lax symmetric monoidal bordism invariant functor $\Catp \to \Sps$. Furthermore, any bordism invariant functor $\Catp \to \E$ valued in a cocomplete $\infty$-category $\E$ is a module over $\Lspace$ in an essentially unique manner.
\item
The functor $\Lspace\colon \Catp \to \Sps$ is Day idempotent and Day convolution with it yields a smashing localisation of $\infty$-operads
\[
\Fun(\Catp,\Sps)^{\otimes} \locadjL{\Lspace \otimes_{\Day} (-)} \Fun^\binv(\Catp,\Sps)^{\otimes} .
\]
\end{enumerate}
\end{corollary}
\begin{proof}
The first part follows from Proposition~\reffour{proposition:universal-unit}
for $\cR=\cR_\binv$ using Propositions~\reffour{prop:closure-properties-of-Rs} and~\reffour{proposition:universal-L-space}. Since arbitrary colimits of equivalences are equivalences, the second part then follows from Proposition~\reffour{proposition:sample-day-smashing}
since $\Lspace$ is small by Remark~\reffour{remark:small-functors}.
\end{proof}









\begin{variant}
Proposition~\reffour{proposition:sample-day-smashing} can also be applied to $\Lspace$ as an $\Einf$-monoid valued functor. For this, consider the collection $\cR^{\oplus-\binv}$ consisting of the diagrams in $\cR^{\binv}$ together with the diagram $\Del^0 = \emptyset^{\triangleleft}\to \Catp$ picking the zero Poincaré category and the diagrams $(\Del^0 \coprod \Del^0)^{\triangleleft} \to \Catp$ witnessing all cartesian products in $\Catp$. Then a functor $\F\colon \Catp \to \Mon_{\Einf}$ is $\cR^{\oplus-\binv}$-local if and only if it is bordism invariant and direct sum preserving. Since $\Catp$ and $\Mon_{\Einf}$ are semi-additivly symmetric monoidal the first two assumptions of Proposition~\reffour{proposition:sample-day-smashing} hold for $\cR^{\oplus-\binv}$ as well. Finally, since $\Lspace$ and $\Poinc$ are both direct sum preserving as $\Mon_{\Einf}$-valued functors and since direct sum preserving functors $\Catp \to \Mon_{\Einf}$ embed in space-valued functors (as the product preserving ones), Proposition~\reffour{proposition:universal-L-space} implies that $\Lspace$ is also initial as a bordism invariant direct sum preserving functor $\Catp \to \Mon_{\Einf}$ under $\Poinc$.
We conclude from Proposition~\reffour{proposition:sample-day-smashing} that $\Lspace$ is idempotent as a $\Mon_{\Einf}$-valued functor and Day convolving with it localises into the subcategory of bordism invariant direct sum preserving functors. A-posteriori, this implies that all bordism invariant direct sum preserving functors to $\Mon_{\Einf}$ are group-like, namely, taking values in $\Grp_{\Einf}$, and that also as an $\Einf$-group valued functor $\Lspace$ is idempotent and localises to bordism invariant direct sum preserving functors.
\end{variant}

\begin{corollary}
\label{corollary:lmonoidal}%
\
\begin{enumerate}
\item
\label{item:L-universal}%
The natural transformation $\SS[\Poinc] \Rightarrow \L$ refines in an essentially unique manner to a lax symmetric monoidal natural transformation among lax symmetric monoidal functors, making $\L$ into the initial 
lax symmetric monoidal bordism invariant and additive functor to $\Sp$. Furthermore, any additive bordism invariant functor $\Catp \to \E$ valued in a stable cocomplete $\infty$-category $\E$ is a module over $\L$ in an essentially unique manner.
\item
\label{item:L-smashing}%
The functor $\L\colon \Catp \to \Spa$ is Day idempotent and Day convolution with it yields a smashing localisation of $\infty$-operads
\[
\Fun(\Catp,\Spa)^{\otimes} \locadjL{\L \otimes_{\Day} (-)} \Funbord(\Catp,\Spa)^{\otimes}.
\]
\end{enumerate}
\end{corollary}
In particular, the bordification $(-)^{\bord}$ of additive functors studied in \S\reftwo{subsection:bordification} satisfies $(-)^\bord \simeq \L \otimes_\Day -$.
\begin{proof}
The first part follows from Proposition~\reffour{proposition:universal-unit}
for $\cR=\cR_\bord$ using Proposition~\reffour{prop:closure-properties-of-Rs} and Theorem~\reftwo{theorem:lisbordgw}. Since arbitrary colimits of equivalences are equivalences and pushouts in spectra are closed under colimits and tensor products, the second part then follows from 
Proposition~\reffour{proposition:sample-day-smashing}
using in addition Remark~\reffour{remark:small-functors}.
\end{proof}









\begin{corollary}
\label{corollary:kl-monoidal}%
The natural transformation $\SS[\Poinc] \Rightarrow \KL$ refines in an essentially unique manner to a lax symmetric monoidal natural transformation among lax symmetric monoidal functors, making $\KL$ into the initial 
lax symmetric monoidal bordism invariant and Karoubi-localising functor to $\Sp$. Furthermore, any bordism invariant Karoubi-localising functor $\Catp \to \E$ valued in a stable cocomplete $\infty$-category $\E$ is a module over $\KL$ in an essentially unique manner.
\end{corollary}
\begin{proof}
We have established in Proposition~\reffour{proposition:universal-KL} above that $\KL$ is the initial bordism invariant and Karoubi-localising functor under $\SS[\Poinc]$, so the result follows from  Proposition~\reffour{proposition:universal-unit}
with $\cR_{\bord} \cup \cR_{\kloc}$.
\end{proof}


\begin{remark}
\label{remark:KL-not-convolving}%
As with $\KGW$ (see Remark~\reffour{remark:not-convolving}), we do not now whether $\KL$ is idempotent, nor whether every functor $G\colon \Catp \to \Spa$ admits a universal Karoubi-localising bordism invariant 
approximation $G \to \mathbb{G}$, but again as for $\KGW$ we do know that such an approximation cannot be obtained by Day convolving with $\KL$, since by Remark~\reffour{rem:contensor-does-not-preserve-KS} functors of the form $\KL(\Funx((\C,\QF),-))$ are generally not Karoubi-localising. 
\end{remark}




\begin{remark}
One can also apply Proposition~\reffour{proposition:sample-day-smashing}
with $\cR = \cR_\binv$ and $\D=\Spa$. We then find that $\SS[\Lspace] \otimes_\Day -$ implements bordification for functors $\Catp \rightarrow \Sp$. Using Corollary \reffour{corollary:GW-universal-multiplicative} above we thus find that $\SS[\Lspace] \otimes_{\Day} \GW$ is the bordification of $\GW$ inside all functors, but retains its additivity. By Theorem~\reftwo{theorem:lisbordgw} we thus find $\L \simeq \SS[\Lspace] \otimes_\Day \GW$.
\end{remark}



	



For completeness sake, we also record the multiplicative properties of K-theory.
\begin{corollary}
\label{corollary:K-universal}%
\
\begin{enumerate}
\item 
The natural transformation $\core \Rightarrow \Kspace$
refines in an essentially unique manner to a lax symmetric monoidal natural transformation among lax symmetric monoidal functors, making $\Kspace$ into the initial lax symmetric monoidal group-like additive functor $\Catx \to \Sps$. Furthermore, any group-like additive functor $\Catx \to \E$ valued in a cocomplete $\infty$-category with finite products $\E$ is a module over $\Kspace$ in an essentially unique manner.
\item
\label{item:k-idempotent}%
The functor $\Kspace\colon \Catp \to \Grp_{\Einf}$ is Day idempotent and Day convolution with it yields a smashing localisation of $\infty$-operads
\[
\Fun(\Catx,\Grp_{\Einf})^{\Day} \locadjL{\Kspace \otimes_{\Day} (-)} \Funadd(\Catx,\Grp_{\Einf})^{\Day} .
\]
\item 
The natural transformation $\SS[\core] \Rightarrow \K$ refines in an essentially unique manner to a lax symmetric monoidal natural transformation among lax symmetric monoidal functors, making $\K$ into the initial lax symmetric monoidal additive functor $\Catx \rightarrow \Spa$. Furthermore, any additive functor $\Catx \to \E$ valued in a stable cocomplete $\infty$-category $\E$ is a module over $\K$ in an essentially unique manner. 
\item
\label{item:K-idempotent}%
The functor $\K\colon \Catp \to \Spa$ is Day idempotent and Day convolution with it yields a smashing localisation of $\infty$-operads
\[
\Fun(\Catx,\Spa)^{\Day} \locadjL{\K \otimes_{\Day} (-)} \Funadd(\Catx,\Spa)^{\Day} .
\]
\item 
The natural transformation $\SS[\core] \Rightarrow \KK$ refines in an essentially unique manner to a lax symmetric monoidal natural transformation among lax symmetric monoidal functors, making $\KK$ into the initial lax symmetric monoidal Karoubi-localising functor $\Catx \rightarrow \Spa$. Furthermore, any Karoubi-localising functor $\Catx \to \E$ valued in a stable cocomplete $\infty$-category $\E$ is a module over $\KK$ in an essentially unique manner.
\end{enumerate}
\end{corollary}
\begin{proof}
We apply Propositions~\reffour{proposition:universal-unit} and \reffour{proposition:sample-day-smashing}
with $\C=\Catx$, target either $\Grp_{\Einf}$ or $\Sp$, and the variants of the collections $\cR_{\add}$ and $\cR_\kloc$ from Definition~\reffour{def:collections-R} obtained by discarding Poincar\'e structures. All the arguments here are the same as their hermitian analogues above, with the exception of Statement~\reffouritem{item:k-idempotent} (whose analogue for $\GWspace$ is false), which requires verifying that the collection of additive functors $\F\colon \Catx \to \Grp_{\Einf}$ is closed under colimits and tensoring with objects in $\Grp_{\Einf}$; indeed, since split Verdier projections admit sections in $\Catx$ we have that $\F$ sends Verdier squares to pullback squares if and only if it sends them to pushout squares.
\end{proof}


\begin{remark}
\label{remark:KK-not-convolving}%
As with $\KGW$ and $\KL$ (see Remarks \reffour{remark:not-convolving} and \reffour{remark:KL-not-convolving}), we do not now whether $\KK$ is idempotent, nor whether every functor $G\colon \Catx \to \Spa$ admits a universal Karoubi-localising approximation $G \to \mathbb{G}$, but as with the previous examples we do know that such an approximation cannot be obtained by Day convolving with $\KK$, since by Remark~\reffour{rem:contensor-does-not-preserve-KS} functors of the form $\KK(\Funx(\C,-))$ are generally not Karoubi-localising, e.g., for $\C = \Dperf(\QQ)$.
As pointed to us by Maxime Ramzi, the same example also shows that Day convolving with $\KK$ cannot localise onto the full subcategory of motivic functors $\Catp \to \Spa$, that is, functors which invert motivic equivalences, or equivalently, the functors which factor through $\Mot$, where we recall from~\cite{ramzi-motivic} that the functor $U\colon \Catx \to \Mot$ exhibits $\Mot$ as the localisation of $\Catx$ by the motivic equivalences. Indeed, we claim that if a functor of the form $\KK(\Funx(\C,-))$ is motivic then $\Om^{\infty}\KK(\Funx(\C,-))$ is Karoubi-localising (which already fails for $\C = \Dperf(\QQ)$). 
To see this, note that if $\KK(\Funx(\C,-))$ is motivic then the space valued functor $\Om^{\infty}\KK(\Funx(\C,-))$ is also motivic. Now since $U$ is symmetric monoidal the operation $\G \mapsto \G(\C\otimes-)$ preserves motivic functors and so its left adjoint $\F \mapsto \F(\Funx(\C,-))$ preserves maps which are local equivalences with respect to motivic functors. Since $\core\Rightarrow \Om^{\infty}\KK = \Map_{\Mot}(\one,U(-))$ is one such local equivalence we conclude that $\core(\Funx(\C,-)) \Rightarrow \Om^{\infty}\KK(\Funx(\C,-))$ is also one, and hence by 2-out-of-3 so is the map
\[
\Om^{\infty}\KK(\Funx(\C,-)) \Rightarrow \Map_{\Mot}(U(\C),U(-)) .
\]
In particular, if the left hand side is motivic then this map would have to be an equivalence, hence the left hand side would be Karoubi-localising, as claimed.
\end{remark}

Using the results of the appendix for more general base $\infty$-operads 
one can deduce further results, for example:

\begin{corollary}
\label{corollary:multgw}%
Let $\E$ be a presentable $\infty$-category and consider the full inclusions
\[
\Funadd(\Catp,\E) \subseteq \Fun(\Catp,\E) \quad ,\quad \Funbord(\Catp,\E) \subseteq \Fun(\Catp,\E)
\]
\[
\Funadd(\Catx,\E) \subseteq \Fun(\Catx,\E) \quad \text{and} \quad \Fun^{\mathrm{binv}}(\Catp,\E) \subseteq \Fun(\Catp,\E).
\]
The the following holds:
\begin{enumerate}
\item
The bottom right inclusion is always reflective with reflection given by $\F \mapsto \Lspace \otimes_{\Day} \F$.
\item
If $\E$ is additive then the bottom left inclusion is reflective with reflection given by $\F \mapsto \Kspace \otimes_{\Day} \F$.
\item 
If $\E$ is stable then the top two inclusions are reflective with reflections given by $\F \mapsto \GW \otimes_{\Day} \F$ and $\F \mapsto \L \otimes_\Day \F$, respectively.
\end{enumerate}
In addition, if we replace $\E$ with a presentably $\cO$-monoidal $\infty$-category for some $\infty$-operad $\cO$ (keeping the additivity or stability assumptions in the relevant cases) then the above reflection functors yield multiplicative localisations on the level of Day convolution $\infty$-operads. When $\cO$ admits a unit object these multiplicative localisations are furthermore smashing.
\end{corollary}


\begin{proof}
Since $\E$ is presentable it carries a canonical colimit preserving action of $\Sps$, and if $\E$ is additive (resp. stable) this action factors through a colimit preserving action of $\Grp_{\Einf}$ (resp. $\Sp$). We may then encode this action via presentably $\MCom$-monoidal $\infty$-category $\M^{\otimes} \to \MCom$, %
to which we can apply Proposition~\reffour{proposition:smashing-day}.
For the last part apply Proposition~\reffour{proposition:multiplicative-day} for the multiplicative statement, given that under the assumptions all functors in question admit local approximations by the first part of the corollary.
As for the smashing statement, let us consider for simplicity the case of $\GW$. Let $\cO^{\otimes}$ be an $\infty$-operad with a unit object $\one_{\cO}$ and $\D^{\otimes} \to \cO^{\otimes}$ a presentably $\cO$-monoidal $\infty$-category, so that $\D_{\one_{\cO}}$ acquires the structure of a presentably symmetric monoidal $\infty$-category. Then post-composition with both the unit functor $u\colon \Sp \to \D_{\one_{\cO}}$ and its right adjoint preserves additive functors and hence $u\GW$ as the initial additive functor under the Day unit $u\Sig^{\infty}_+\Poinc\colon \Catp \to \D_{\one_{\cO}}$. In addition, if $a \in \cO$ is any object then $\D_a$ is a stable presentable $\infty$-category and the action of $\Sp$ on $\D_a$ factors through an action of $\D_{\one_{\cO}}$, 
so that the operations $u\GW\otimes_{\Day}(-)$ and $\GW\otimes_{\Day} (-)$ on $\D_a$ coincide. The desired result now follows from Proposition~\reffour{proposition:smashing-day}. The argument in the case of $\L,\Lspace$ and $\Kspace$ are the same.
\end{proof}







Having established the multiplicative properties of the K, GW, and L-functors, let us finally spell out their interaction. To this end, recall from Corollary~\refone{corollary:hyp-is-adjoint} the adjunctions  
\begin{equation}
\label{equation:two-sided}%
\begin{tikzcd}
	\Catx \ar[r, bend right,"\Hyp"'] \ar[r, bend left,"\Hyp"] & \Catp \ar[l, "\mathrm{fgt}"']
\end{tikzcd}
\end{equation}
exhibiting $\Hyp$ as both left and right adjoint to the forgetful functor $\fgt$. Since the functors $\Hyp$ and $\fgt$ preserves zero objects and split (Poincaré-)Verdier squares, restriction along $\Hyp$ and $\fgt$ is an operation which preserves additive functors. We then obtain an induced two sided adjunction 
\[
\begin{tikzcd}
	\Funadd(\Catp,\Spa) \ar[r, "\Hyp^*"]  & \Funadd(\Catx,\Spa) \ar[l, bend right,"\fgt^*"'] \ar[l, bend left,"\fgt^*"]
\end{tikzcd}
\]
on the level of additive functors.
In Proposition~\reftwo{proposition:bordrecoll}, it was shown that this two-sided adjunction yields a stable recollement
\[
\begin{tikzcd}
[column sep=7ex]
\Funbord(\Catp,\Spa) \ar[r] & 
\Funadd(\Catp,\Spa) 
\ar[r] 
\ar[l,bend left=25,shift left=1.5ex,start anchor=west,end anchor=east, "(-)^{\rbord}"] 
\ar[l,bend right=25,shift right=1.5ex,start anchor=west,end anchor=east,"(-)^{\bord}"'] 
\ar[l,phantom,shift left=1.2ex,start anchor=west,end anchor=east,"\perp"] 
\ar[l,phantom,shift right=1.3ex,start anchor=west,end anchor=east,"\perp"]
& \Funadd(\Catx,\Spa)^\hC \ .
\ar[l,bend left=25,shift left=1.5ex,start anchor=west,end anchor=east,"{(- \circ \fgt)^\hC}"] 
\ar[l,bend right=25,shift right=1.5ex,start anchor=west,end anchor=east,"{(- \circ \fgt)_\hC}"']
\ar[l,phantom,shift left=1.2ex,start anchor=west,end anchor=east,"\perp"] 
\ar[l,phantom,shift right=1.3ex,start anchor=west,end anchor=east,"\perp"]
\end{tikzcd}
\]
In~\cite{shah2021recollements}*{Definition 2.20} Shah defines the notion of a multiplicative stable recollement with respect to a symmetric monoidal structure on the ambient $\infty$-category, by requiring that the left Bousfield localisations to the two constituents of the recollement are compatible with said structure. 
His definition readily extends to the case where the ambient $\infty$-category is not a symmetric monoidal $\infty$-category but only an $\infty$-operad, by requiring these left Bousfield localisations to be multiplicative in the sense of Definition~\reffour{definition:multiplicative-localisation}. 
Returning to the above recollement, we may view the ambient term as the $\infty$-category of colours underlying the Day convolution full suboperad $\Funadd(\Catp,\Spa)^{\Day} \subseteq \Fun(\Catp,\Spa)^{\Day}$. The two constituents of the recollement can then be viewed as full suboperads of $\Funadd(\Catp,\Spa)^{\Day}$. For the left constituent this just gives the Day convolution full suboperad $\Funbord(\Catp,\Spa)^{\Day} \subseteq \Fun(\Catp,\Sp)^{\Day}$. As for the right constituent, we have the following:

\begin{lemma}
\label{lemma:sanity-check}%
The full suboperad of $\Funadd(\Catp,\Spa)^{\Day}$ spanned by the image of the fully-faithful right adjoint 
\[
(- \circ \fgt)^\hC \colon \Funadd(\Catx,\Spa)^{\hC} \to \Funadd(\Catp,\Spa)
\]
naturally coincides with the homotopy fixed points $(\Funadd(\Catx,\Spa)^{\Day})^{\hC}$ of the Day convolution full suboperad $\Funadd(\Catx,\Spa)^{\Day} \subseteq \Fun(\Catx,\Spa)^{\Day}$. 
\end{lemma}
\begin{proof}
Recall that the forgetful functor $\fgt\colon \Catp \to \Catx$ is symmetric monoidal and $\Ct$-equivariant, where $\Ct$ acts trivially on $\Catp$. Pre-composition with $\fgt$ hence determines a $\Ct$-equivariant map of $\infty$-operads
\[
(-\circ \fgt)\colon \Fun(\Catx,\Spa)^{\Day} \to \Fun(\Catp,\Spa)^{\Day} ,
\]
with $\Ct$ acting trivially on the target.
Passing to homotopy fixed points and using that $\Fun(\Catp,-)^{\Day}$ is a right adjoint and hence limit preserving we obtain a composite map of $\infty$-operads
\[
(\Fun(\Catx,\Spa)^{\Day})^{\hC} \to (\Fun(\Catp,\Spa)^{\Day})^{\hC}  = \Fun(\Catp,\Spa^{\B\Ct})^{\Day} \xrightarrow{(-)^{\hC}\circ (-)} \Fun(\Catp,\Spa)^{\Day} ,
\]
where in the last step we composed with the $\infty$-operad map induced on Day convolution $\infty$-operads by the lax symmetric monoidal functor $(-)^{\hC}\colon \Spa^{\B\Ct} \to \Spa$ (viewed as an $\infty$-operad map). Now the full suboperad inclusion $\Funadd(\Catx,\Spa)^{\Day} \subseteq \Fun(\Catx,\Spa)^{\Day}$ induces a full suboperad inclusion after taking homotopy fixed points, and the $\infty$-operad map constructed above sends this full suboperad to the full suboperad $\Funadd(\Catp,\Spa)$, thus giving an $\infty$-operad map
\[
(-\circ \fgt)^{\hC}\colon (\Funadd(\Catx,\Spa)^{\Day})^{\hC} \to \Funadd(\Catp,\Spa),
\]
refining the fully-faithful functor $(-\circ \fgt)^{\hC}\colon \Funadd(\Catx,\Spa)^{\hC} \hrar \Funadd(\Catp,\Spa)$ on the level of colours. To finish the proof it will hence suffice to show that $(-\circ \fgt)^{\hC}$ is also fully-faithful as a map of $\infty$-operads, that is, for every sequence of $\Ct$-equivariant additive functors $\F_1,...,\F_n,\G\colon \Catx \to \Spa$ (with respect to the $\mop$ action on $\Catx$ and the trivial action on $\Spa$) the induced map
\[
\Nat(\otimes_n \circ (\F_1 \times ... \times \F_n),\G \circ \otimes_n)^{\hC} \to \Nat([\otimes_n \circ ([\F_1 \circ \fgt]^{\hC} \times ... \times [\F_n \circ \fgt]^{\hC}),[\G \circ \fgt]^{\hC} \circ \otimes_n)
\]
is an equivalence, where we write $\otimes_n$ for the $n$-fold tensor product of either $\Catx,\Catp$ or $\Spa$. 
Now since $\fgt$ is symmetric monoidal we can rewrite this map as
\[
\Nat(\otimes_n \circ (\F_1 \times ... \times \F_n),\G \circ \otimes_n)^{\hC} \to \Nat([\otimes_n \circ ([\F_1 \circ \fgt]^{\hC} \times ... \times [\F_n \circ \fgt]^{\hC}),[\G \circ \otimes_n\circ (\fgt \times ... \times \fgt)]^{\hC}) .
\]
Next, since $\Hyp$ is in particular right adjoint to $\fgt$, the functor 
\[
[(-) \circ (\fgt \times ... \times \fgt)]^{\hC}\colon \Fun(\Catx \times ... \times \Catx,\Spa)^{\hC} \to \Fun(\Catp \times ... \times \Catp,\Spa)
\]
has a left adjoint given by
\[
(-) \circ (\Hyp \times ... \times \Hyp)\colon \Fun(\Catp \times ... \times \Catp,\Spa) \to \Fun(\Catx \times ... \times \Catx,\Spa)^{\hC}.
\]
The map that we want to show is an equivalence can hence be identified via adjunction with the map
\[
\Nat(\otimes_n \circ (\F_1 \times ... \times \F_n),\G \circ \otimes_n)^{\hC} \to \Nat([\otimes_n \circ ([\F_1 \circ \fgt \circ \Hyp]^{\hC} \times ... \times [\F_n \circ \fgt \circ \Hyp]^{\hC}),\G \circ \otimes_n)^{\hC}
\]
induced by pre-composition with the composite 
\[
\otimes_n\circ ([\F_1 \circ \fgt\circ \Hyp]^{\hC},...,[\F_n \circ \fgt \circ \Hyp]^{\hC}) \Rightarrow [\otimes_n\circ (\F_1 \circ \fgt\circ \Hyp,...,\F_n \circ \fgt \circ \Hyp)]^{\hC} \Rightarrow   \otimes_n\circ (\F_1,...,\F_n) ,
\]
where the second map is the counit of the above adjunction. But this composite is simply the map induced after $\otimes_n \circ$ by the tuple of counits $[\F_i \circ \fgt \circ \Hyp]^{\hC} \Rightarrow \F_i$
of the adjunction $(-)\circ \Hyp\dashv [(-)\circ \fgt]^{\hC}$ for $i=1,...,n$. The desired result now follows from the fact that the counit $[\F \circ \fgt \circ \Hyp]^{\hC} \Rightarrow \F$ is an equivalence whenever $\F \colon \Catx \to \Spa$ is an additive $\Ct$-equivariant functor.
\end{proof}




\begin{proposition}
The above stable recollement is multiplicative with respect to the Day convolution $\infty$-operad structure on $\Funadd(\Catp,\Spa)$, yielding two left Bousfield localisations of $\infty$-operads
\[
\begin{tikzcd}
[column sep=7ex]
\Funbord(\Catp,\Spa)^{\Day} \ar[r] & 
\Funadd(\Catp,\Spa)^{\Day} 
\ar[r]  
\ar[l,bend right=25,shift right=1.5ex,start anchor=west,end anchor=east,"(-)^{\bord}"'] 
\ar[l,phantom,shift right=1.3ex,start anchor=west,end anchor=east,"\perp"]
& (\Funadd(\Catx,\Spa)^{\Day})^\hC \ .
\ar[l,bend left=25,shift left=1.5ex,start anchor=west,end anchor=east,"{(- \circ \fgt)^\hC}"] 
\ar[l,phantom,shift left=1.2ex,start anchor=west,end anchor=east,"\perp"] 
\end{tikzcd}
\]
\end{proposition}

\begin{proof}
We need to verify that the left localisations 
\begin{enumerate}
\item[-] $(-)^{\brd}\colon \Funadd(\Catp,\Spa) \to \Funbord(\Catp,\Spa)$ and,
\item[-] $\Hyp^*\colon \Funadd(\Catp,\Spa) \to \Funadd(\Catx,\Spa)^{\hC}$
\end{enumerate}
are multiplicative with respect to Day convolution. 
For $(-)^{\brd}$ this follows from Corollary~\reffour{corollary:lmonoidal}\;\reffouritem{item:L-smashing} and Remark~\reffour{remark:full-suboperad}.
For $\Hyp^*$, it will suffice by Lemma~\reffour{lemma:canonical-colimit} to show that $\Hyp^*$-equivalences are closed under Day convolution with (small) functors of the form $E \otimes \SS[\Map_{\Catp}((\C,\QF),-)]\colon \Catp \to \Spa$. Concretely, we need to for $\F\Rightarrow \F'$ a natural transformation of functors $\Catp \to \Spa$ such that $\F(\Hyp(\D)) \to \F'(\Hyp(\D))$ is an equivalence for every $\D \in \Catx$ also the map $E \otimes \F(\Funx((\C,\QF),\Hyp(\D))) \to E \otimes \F'(\Funx((\C,\QF),\Hyp(\D)))$ is an equivalence, which follows from the natural equivalence 
$\Funx((\C,\QF),\Hyp(-)) \simeq \Hyp(\Funx(\C,-))$.
\end{proof}

Applying~\cite{shah2021recollements}*{Corollary 2.29}, which applies verbatim to $\infty$-operads rather than symmetric monoidal $\infty$-categories, we now obtain the following symmetric monoidal refinement of the structure square for additive functors:

\begin{corollary}
\label{corollary:generic-fracture-monoidal}%
Associating to a functor $\F \colon \Catp \to \Spa$ 
its fundamental fibre square
\[
\begin{tikzcd}
	\F \ar[r] \ar[d] & \F^\brd \ar[d]   \\
	(\F^\hyp)^{\hC} \ar[r] & (\F^\hyp)^{\tC} 
\end{tikzcd}
\]
naturally refines to a map of $\infty$-operads $\Fun(\Catp,\Spa)^\Day \to \Fun(\Delta^1 \times \Delta^1,\Fun(\Catp,\Spa)^\Day)$. In particular, the fundamental fibre square of $\F$ naturally refines to a square of lax symmetric monoidal functors once $\F$ refines to a lax symmetric monoidal functor.
\end{corollary}



\begin{corollary}
\label{corollary:kgw-kl-monoidal}%
\label{corollary:GW-symmetric-monoidal}%
The entire cube
\[
\begin{tikzcd}
[column sep=0.5em, row sep=0.7em]
  & \KGW\ar[rr]\ar[dd] && \KL\ar[dd]\\
\GW\ar[ru]\ar[rr, crossing over]\ar[dd] && \L\ar[ur]\\
  & \KK^\hC\ar[rr] && \KK^\tC\\
\K^\hC\ar[rr] \ar[ru] && \K^\tC\ar[uu,<-, crossing over] \ar[ru]\\
\end{tikzcd}
\]
canonically refines to a diagram of lax symmetric monoidal functors $\Catp \rightarrow \Sp$. Here, the lax symmetric monoidals on the functors in the top face are those of Corollaries~\reffour{corollary:kl-monoidal}, 
\reffour{corollary:lmonoidal}\;\reffouritem{item:L-universal},
\reffour{corollary:kgw-multiplicative} and 
\reffour{corollary:GW-symmetric-monoidal}\;\reffouritem{item:GW-universal}, 
while those on the functors in the bottom face are induced by the lax symmetric monoidal structures of Corollary~\reffour{corollary:K-universal}. %
\end{corollary}
\begin{proof}
Note first that the lax symmetric monoidal structures induced on $\L = \GW^\mathrm{bord}$ and $\KL = \KGW^\mathrm{bord}$ by Corollary~\reffour{corollary:generic-fracture-monoidal} 
necessarily agree with those of Corollaries \reffour{corollary:lmonoidal}
and \reffour{corollary:kl-monoidal} respectively, on account of the uniqueness property of the latter. %
Secondly, the map $\GW \Rightarrow \KGW$ canonically refines to one of lax symmetric monoidal functors, since $\GW$ is initial both as an additive functor under $\SS[\Poinc]$ (Corollary~\reftwo{corollary:univ-GW-spectrum}) and as an additive lax symmetric monoidal functor (Corollary~\reffour{corollary:GW-universal-multiplicative}), while at the same time $\SS[\Poinc]$ is the initial lax symmetric monoidal functor $\Catp \to \Sp$ and the composite $\SS[\Poinc] \Rightarrow \GW \Rightarrow \KGW$ is a map of lax symmetric monoidal functors (Corollary~\reffour{corollary:kgw-multiplicative}). We may then apply Corollary~\reffour{corollary:generic-fracture-monoidal} to the lax monoidal transformation $\GW \to \KGW$.
\end{proof}
Using the fact that the pullback diagram
\[
\begin{tikzcd}
\Sp^\gCt \ar[rr,"{(-)^{\varphi\Ct} \Rightarrow (-)^\tC}"] \ar[d,"\mathrm{fgt}"] && \Ar(\Sp) \ar[d,"t"]\\
\Sp^\hC \ar[rr,"{(-)^\tC}"] && \Sp,
\end{tikzcd}
\]
describing genuine $\Ct$-spectra as Tate-squares as above is one of $\infty$-operads, Corollary~\reffour{corollary:kgw-kl-monoidal} can equivalently be stated as:

\begin{corollary}
\label{corollary:real-is-lax-symmetric-monoidal}%
The natural transformation $\KR \Rightarrow \KKR \colon \Catp \rightarrow \Sp^\gCt$ canonically refines to one of lax symmetric monoidal functors.
\end{corollary}


This result in particular provides an $\Einf$-structure to the genuine $\Ct$-spectra $\KR(\C,\QF)$ and $\KKR(\C,\QF)$ whenever $(\C,\QF)$ is a symmetric monoidal Poincar\'e $\infty$-category. In this paper, we will not discuss the more subtle structure of a normed $\Einf$-algebra structure (sometimes also called a $\Ct$-$\Einf$-structure) on $\KR(\C,\QF)$. 









We finish this subsection by explaining how to deduce from the above results multiplicative statements about $\GW$ (or any of the other functors considered above) of rings with form parameters. More precisely, let $(\Ringp)^{\otimes}$ be the symmetric monoidal $\infty$-category sitting in the pullback square
\[
\begin{tikzcd}[column sep=40pt]
(\Ringp)^{\otimes}\ar[r]\ar[d] & (\Catpi)^{\otimes}\ar[d] \\
\Alg_{\Eone}(\Sp)^{\otimes} \ar[r,"\Modp{-}"] & (\Catxi)^{\otimes}
\end{tikzcd}
\]
of symmetric monoidal functors, where $\Catxi$ is the symmetric monoidal localisation of $\Catx$ by the Karoubi equivalences, and similarly for $\Catpi$. A similar construction using $\Alg_{\Einf}$ in place of $\Alg_{\Eone}$ was considered in~\cite{poincare-rings}, who coined the term ``Poincaré rings''.
By the classification of Poincaré structures on module $\infty$-categories (see \S\refone{subsection:classification-in-modules}) we may identify objects in $\Ringp$ with tuples $(A,M,N,\alp\colon N \to M^{\tC})$ where $A$ is an $\Eone$-ring spectrum and $\M := (M,N,\alp)$ is a module with genuine involution over $A$, or, equivalently, an $\N A$-module in the $\infty$-category $\Sp^{\gCt}$ of genuine $\Ct$-spectra, where $\N A \in \CAlg(\Sp^{\g\Ct})$ is the Hill-Hopkins-Ravanel norm of $A$ (whose underlying spectrum is $A \otimes A$ and whose geometric fixed points is $A$), see Remark~\refone{remark:NormModules}. %
The symmetric monoidal product on $(\Ringp)^{\otimes}$ is then given by $(A,\M) \otimes (A',\M') = (A \otimes A',\M \otimes \M')$, where $\M \otimes \M'$ is the tensor product in genuine $\Ct$-spectra, equipped with its induced $A \otimes A'$-module structure, see Example~\refone{example:tensor-of-rings}. Now recall that the $\infty$-category $\Sp^{\gCt}$ admits a t-structure, which we will refer to as the Mackey t-structure, where an $E \in \Sp^{\gCt}$ is Mackey connective (resp. Mackey coconnective) if and only if the underlying spectrum $\bar{E} \in \Sp$ and the genuine fixed points $E^{\Ct} \in \Sp$ are connective (resp. coconnective). In other words, this is the pointwise t-structure after identifying $\Sp^{\gCt}$ with spectral Mackey functors, and its heart consists exactly of the $\Ab$-valued Mackey functors. It is also convenient to note that connectivity (though not coconnectivity) can be tested on underlying objects and geometric fixed points. This t-structure is multiplicative in the sense that the monoidal unit is Mackey connective and the collection of Mackey connective objects is closed under tensor product. In addition, if $A$ is a connective ring spectrum then $\N A$ is Mackey connective.
Finally, let us say that a module with involution over $A$ is projective if it is projective 
as an $A$-modules (that is, is a retract of $A^n$ for some $n$) with respect to either of its $A$-module structures, and that an $\N A$-module is projective if the underlying module with involution is projective in the above sense. We recall from Remark~\refone{remark:form-parameters-and-modules-with-genuine-involution} that when $A$ is a discrete ring the data of a Mackey discrete $\N A$-module corresponds exactly to that of a form parameter over $A$. We say that a form parameter $\lambda$ on $A$ is Poincaré if the associated module with involution $M$ over $A$ is projective and invertible, which is equivalent to saying that 
the hermitian structure $\QF^{\g\lambda}_M$ on $\Dperf(A)$ is Poincaré and its duality preserves $\Proj(R)$ (see~\cite{comparison}*{Example 3.2.4(4)} for a discussion, though note the notation clash with loc.\ cit.\ where the term invertible is used for what we call here Poincaré). The $\infty$-category $\FPMp$ of Poincaré form parameters can then be identified with the full subcategory of $\Ringp$ spanned by those $(A,\M)$ such that $A$ is discrete and $\M$ is projective and Mackey discrete. It sits inside the intermediate full subcategory 
\[
\FPMp \subseteq \Ringp_{\geq 0} \subseteq \Ringp
\]
where $\Ringp_{\geq 0}$ is spanned by those $(A,\M)$ such that $A$ is connective and $\M$ is projective and Mackey connective.
In particular, since the Mackey t-structure is multiplicative we have that $\Ringp_{\geq 0}$ contains the unit $(\SS,\SS^{\uni})$ of $\Ringp$ and is closed under tensor products, so that it inherits from $\Ringp$ a symmetric monoidal structure.


\begin{proposition}%
\label{proposition:sym_mon_structure_on_cat_of_form_parameters}%
The full subcategory $\FPMp \subseteq \Ringp_{\geq 0}$ is a symmetric monoidal localisation of $\Ringp_{\geq 0}$ with reflection functor $\Ringp_{\geq 0} \to \FPMp$ given by $(A,\M) \mapsto (\pi_0(A),\pi^{\mack}_0(\M))$, where $\pi^{\mack}_0$ denotes the $\pi_0$ functor with respect to the Mackey t-structure. In particular, $\FPMp$ inherits a symmeric monoidal structure making its inclusion in $\Ringp_{\geq 0}$ lax symmetric monoidal. The unit of this symmetric monoidal structure is $(\ZZ,\uline{\ZZ}^{\burn})$, where $\uline{\ZZ}^{\burn}$ is the Burside $\N\ZZ$-module, whose underlying $(\ZZ \otimes_{\SS} \ZZ)$-module is $\ZZ$ and whose genuine fixed points object is the Burnside ring of $\ZZ$, see Example~\refthree{example:dotto-ogle}. 
\end{proposition}

Concretely, the induced tensor product in $\FPMp$ is given by 
\[
(R,\M)\otimes (R',\M) = (\pi_0(R \otimes_{\SS} R'),\pi^{\mack}_0(\M\otimes\M')) = (R\otimes_{\ZZ} R',\M\otimes^{\mack} \M')
\]
where $\otimes^{\mack}$ is the induced tensor product on the heart of the Mackey t-structure, also known as the box product of $\Ab$-valued Mackey functors, see Definition~\reffour{definition:box-mackey} for an explicit algebraic description. We note that the inclusion $\Catpi \subseteq \Catp$ as idempotent complete Poincaré $\infty$-categories is lax symmetric monoidal, and so we get that the composite functor
\[
\FPMp \to \Catp \quad\quad (R,\lambda) \mapsto (\Dperf(R),\QF^{\g\lambda})
\]
is lax symmetric monoidal.
In particular, $\GW$ and $\L$ (and similarly all the functors considered above) promote to lax symmetric monoidal functors on $\FPMp$, and for any $(R, \lambda)$ in $\FPMp$, its $\GW$ and $\L$-theories are canonically modules over $\GW^{\burn}(\ZZ)$ and $\L^{\burn}(\ZZ)$ respectively. 
 
 
\begin{proof}[Proof of Proposition~\reffour{proposition:sym_mon_structure_on_cat_of_form_parameters}]
Let $(A,\M) \in \Ringp_{\geq 0}$, so that $\pi_0(A)$ inherits the structure of a ring. We then observe that the induced map $\N A \to \N\pi_0(A)$ in $\Sp^{\gCt}$ has a Mackey 1-connective fibre, and so induces an isomorphism on $\pi_0^{\mack}$. Now since the Mackey t-structure is multiplicative we may refine $\pi_0^{\mack}$ to a left adjoint to the inclusion of Mackey discrete algebra objects inside Mackey connective algebra objects. It hence follows that the algebra map $\N A \to \pi_0^{\mack}(\N A)$ factors uniquely through an algebra map $\N\pi_0(A) \to \pi^{\mack}_0(\N A)$, through which we may view $\pi_0^{\mack}(\M)$ as an $\N\pi_0(A)$-module, or, equivalently, as a module with genuine involution over $\pi_0(A)$. Now since the Mackey t-structure is pointwise under the Mackey functor description we have that the underlying $\pi_0(A)\otimes_{\SS} \pi_0(A)$-module of $\pi_0^{\mack}(\M)$ is given by $\pi_0(M)$, where $M$ is the $A \otimes A$-module underlying $\M$. Since $M$ is projective as an $A$-module we have that $\pi_0(M)$ is projective as a $\pi_0(A)$-module.
We now claim that $\pi_0(M)$ is invertible as a module with involution over $\pi_0(A)$, that is, that the induced map
\[
\pi_0(A) \to \map_{\pi_0(A)}(\pi_0(M),\pi_0(M))
\]
is an equivalence. %
For this, note that the left hand side is a-priori connective since $\pi_0(M)$ is projective. Refining $\pi_0$ to a left adjoint to the inclusion $\Mod_{\pi_0(A)}(\Ab) \to \Mod_{A}(\Sp)$ we may then replace the right hand side with $\map_{A}(M,\pi_0(M))$. We then note that $\map_A(M,-)$ commutes with truncation since $M$ is projective and connective, so that the desired claim follows from the fact that the original $M$ was invertible as a module with involution over $A$.
We now claim that the resulting map
\[
(A,\M) \to (\pi_0(A),\pi_0(\M))
\]
is Poincaré, that is, that the induced map $M\otimes_A \pi_0(A) \to \pi_0(M)$ is an equivalence. Indeed, the fact that the standard t-structure is multiplicative implies that this map is an isomorphism on $\pi_0$, where on the other hand the both sides are discrete since $M$ is projective as an $A$-module.

Finally, we now show that $(\pi_0(A),\pi_0(\M))$ is initial in $\FPMp_{(A,\M)/}$.  Unwinding all definitions and using $\pi_0$ and $\pi_0^{\mack}$ are left adjoints to the inclusion of discrete objects inside connective one, it will suffice to show that if $(R,\mathcal{N})$ is an object of $\FPMp$ equipped with a ring map $\pi_0(A) \to R$ and an $\N\pi_0(A)$-module map $\pi_0^{\mack}(\M) \to \mathcal{N}$ then the map $(\pi_0(A),\pi_0(\M)) \to (R,\mathcal{N})$ is Poincaré as soon as the composite $(A,\M) \to (\pi_0(A),\pi^{\mack}_0(\M)) \to (R,\mathcal{N})$ is Poincaré. In other words, the induced map $\pi_0(M)\otimes_{\pi_0(A)} R \to N$ is an equivalence as soon as the induced map $M \otimes_{A} R \to N$ is so, where $N$ is the underlying module with involution over $R$. Indeed, this is simply because the map $M\otimes_A \pi_0(A) \to \pi_0(M)$ is an equivalence, as we showed above.

We have thus shown that the inclusion $\FPMp \subseteq \Ringp_{\geq 0}$ is reflective with reflection functor $F(A,M) = (\pi_0(A),\pi^{\mack}_0(\M))$. To show that the resulting localisation is multiplicative we need to show the collection of $F$-equivalences are closed under taking tensor products with fixed objects. Now by construction the $F$-equivalence are exactly the maps $(A,\M) \to (B,\mathcal{N})$ such that the ring spectrum map $A \to B$ is a $\pi_0$-isomorphism and the $\N A$-module map $\M \to \mathcal{N}$ is a $\pi_0^{\mack}$-isomorphism. This collection of maps is closed under tensoring with fixed objects simply because the standard t-structure on $\Sp$ and the Mackey t-structure on $\Sp^{\g\Ct}$ are both multplicative.
\end{proof}












\subsection{Examples of multiplicative structures on Grothendieck-Witt- and L-spectra}
\label{subsection:examples}%

In this section we will collect some examples of naturally occurring symmetric monoidal Poincaré $\infty$-categories (and modules over them) to which the multiplicative invariants summarised in the cube of Corollary~\reffour{corollary:kgw-kl-monoidal} can be applied, yielding cubes of commutative ring spectra (and modules over such), and we will describe the multiplicative structure on several (previously additively computed) Grothendieck-Witt and L-spectra. We begin with categorical examples.



\begin{example}
\label{example:initial-terminal}%
If $(A,\sig)$ is an $\Einf$-ring spectrum with involution, the $\infty$-category $\Modp{A}$ of compact $A$-modules carries a symmetric monoidal structure and $\Dual_{\sig}$ is a symmetric monoidal duality on $\Modp{A}$. The associated Tate construction $A^{\tC}$ then inherits the structure of an $\Einf$-$A$-algebra spectrum via the twisted Tate-Frobenius and symmetric monoidal Poincaré refinements of $(\Dperf(A),\Dual_{\sig})$ are classified by $\Einf$-$A$-algebras $N$ (representing the linear part of the associated Poincaré refinement) equipped with an $\Einf$-$A$-algebra map to $A^{\tC}$. Among those, the initial such algebra is given by the unit map $A \to A^{\tC}$, and corresponds to the Tate Poincaré structure $\QF^{\tate}_{\sig}$, while the terminal such algebra is given by the identity map $A^{\tC} \to A^{\tC}$, and corresponds to the symmetric Poincaré structure $\QF^{\sym}_{\sig}$. For $A=\SS$ the Tate Poincaré structure on $\Modp{\SS} = \Spaf$ coincides with the universal one of Example~\reffour{example:universal}.
In particular, for $(A,\sig)$ as above, we have a (pullback) squares of $\Einf$-ring spectra 
\[
\begin{tikzcd}
\GW^{\tate}(A,\sig) \ar[r]\ar[d] & \L^{\tate}(A,\sig) \ar[d] &\KGW^{\tate}(A,\sig) \ar[r]\ar[d] & \KL^{\tate}(A,\sig) \ar[d] \\
\GW^{\sym}(A,\sig)\ar[r] & \L^{\sym}(A,\sig) & \KGW^{\sym}(A,\sig)\ar[r] & \KL^{\sym}(A,\sig)..
\end{tikzcd}
\]
In addition, if $B$ is an $\Eone$-$A$-algebra equipped with a module with involution $M$ whose $\Ct$-action is $A$-sesquilinear then $(\Modp{B},\QF^{\sym}_M)$ is a module over $(\Modp{A},\QF^{\sym}_{\sig})$ and any Poincaré refinement $(\Modp{B},\QF_M)$ of $\Dual_{M}$ is a module over $(\Modp{A},\QF^{\tate}_{\sig})$. Consequently, $\GW^\s(\Modp{B},\Dual_M))$ and $\L^\s(\Modp{B},\Dual_M))$ are modules over $\GW^s(A,\sig)$ and $\L^\s(A,\sig)$, respectively, and for any Poincar\'e refinement $\QF_M$ of $\Dual_M$, we find that $\GW(\Modp{B},\QF_M)$ and $\L(\Modp{B},\QF_M)$ are modules over $\GW^\tate(A,\sig)$ and $\L^\tate(A,\sig)$, respectively. The same holds for $\KGW$ and $\KL$ in place of $\GW$ and $\L$.
\end{example}


\begin{example}
\label{example:geq-zero}%
If $(A,\sig)$ is a connective $\Einf$-ring spectrum with involution, the Poincaré structure $\QF^{\geq 0}_{\sig}$ on $\Modp{A}$ is symmetric monoidal: its associated $\Einf$-$A$-algebra is the connective cover $\tau_{\geq 0} A^{\tC}$ of $A^{\tC}$, which is indeed an $\Einf$-$A$-algebra since the t-structure on $\Mod_A$ is compatible with its symmetric monoidal structure. In fact, $\QF^{\geq0}_{\sig}$ constitutes the terminal symmetric monoidal Poincaré refinement of $(\Modp{A},\Dual_{\sig})$ with connective linear part. If $B$ is a connective $\Eone$-$A$-algebra equipped with an $A$-sesquilinear module with involution $M$ then $(\Modp{B},\QF^{\geq m}_M)$ is a module over $(\Modp{A},\QF^{\geq 0}_{\sig})$ for all $m\in \ZZ \cup \{ \pm \infty \}$. Consequently, $\GW(\Modp{B},\QF^{\geq m}_M)$ and $\L(\Modp{B},\QF^{\geq m}_M)$ are modules over $\GW(\Modp{A},\QF^{\geq0}_\sig)$ and $\L(\Modp{A},\QF^{\geq0}_\sig)$, respectively. The same holds for $\KGW$ and $\KL$ in place of $\GW$ and $\L$.
\end{example}


\begin{example}
\label{example:ideal}%
If $(A,\sig)$ is an $\Einf$-ring spectrum with involution, then $(\Modp{A},\QF^\qdr_\sig)$ is not a symmetric monoidal Poincar\'e $\infty$-category, but it is a module over $(\Modp{A},\QF^\sym_\sig)$. In fact, the canonical functor $(\Modp{A},\QF^\qdr_\sig) \to (\Modp{A},\QF^\sym_\sig)$ is a Smith ideal, that is, a commutative algebra in $\Fun([1],\Catp)$ with respect to the Day convolution symmetric monoidal structure for the cartesian product on $[1]$, see \S\reffour{subsection:normal-L} for details. Consequently, $\GW^\qdr(A,\sig)$ and $\L^\qdr(A,\sig)$ are Smith ideals in $\GW^\s(A,\sig)$ and $\L^\s(A,\sig)$, respectively, whose quotients agree (by the fundamental fibre sequence) and are called symmetric normal $\L$-theory; 
see again \S\reffour{subsection:normal-L} for more on normal L-theory. If $A$ is connective and $m\geq 0$, then $(\Modp{A},\QF^{\geq m}_\sig) \to (\Modp{A},\QF^{\geq0}_\sig)$ is a Smith ideal, so $\GW^{\geq m}(A,\sig)$ and $\L^{\geq m}(A,\sig)$ are Smith ideals in $\GW^{\geq 0}(A,\sig)$ and $\L^{\geq0}(A,\sig)$, respectively.
The same holds for $\KGW$ and $\KL$ in place of $\GW$ and $\L$.
\end{example} 






\begin{example}
\label{example:modules-over-genuine}%
If $(R,\sig)$ is a discrete commutative ring with involution then we may also consider the genuine Frobenius Poincar\'e structure $\QF^{\g\fr}_\sig$ of $\Dperf(R)$. It participates in the following sequence of symmetric monoidal refinements of $(\Dperf(R),\Dual_{\sig})$
\[
\QF^{\tate}_{\sig} \Rightarrow \QF^{\g\fr}_{\sig} \Rightarrow \QF^{\gs}_{\sig} \Rightarrow \QF^{\sym}_{\sig} ,
\]
where $\QF^{\tate}_{\sig}$ and $\QF^{\sym}_{\sig}$ are the initial and terminal refinements as in Example~\reffour{example:initial-terminal} and the genuine symmetric structure $\QF^{\gs}_{\sig} = \QF^{\geq 0}_{\sig}$ is the terminal refinement with connective linear part as in of Example~\reffour{example:geq-zero}.
The genuine Frobenius Poincaré structure $\QF^{\g\fr}_{\sig}$ on $\Dperf(R)$ is obtained from $\QF^{\tate}_{\sig}$ by applying the reflection functor of Proposition~\reffour{proposition:sym_mon_structure_on_cat_of_form_parameters}, and 
constitutes the initial symmetric monoidal %
Poincaré structure refining $\Dual_{\sig}$ whose associated module with involution $(R,N^{\g\fr} \to R^{\tC})$ is Mackey 0-truncated. %
The $\Einf$-$R$-algebra spectrum $N^{\g\fr}$ is connective with $\pi_n(N^{\g\fr}) = \pi_n(R^{\tC})$ for $n \geq 2$, $\pi_1(N^{\g\fr}) = 0$ and $\pi_0(N^{\g\fr}) = R$, so that $(R,N^{\g\fr})$ is also Mackey connective and $\QF^{\g\fr}_{\sigma}$ corresponds to the Frobenius form parameter of
Example~\reffour{example:form_parameter_over_Tate_Poincare}.
For $R=\ZZ$ the genuine Frobenius structure $\QF^{\g\fr}_{\ZZ}$ coincides with the Burnside Poincaré structure %
appearing in Example~\refthree{example:dotto-ogle}.

If $(R,\sig)$ is a discrete commutative ring with involution and $B$ is a discrete associative $R$-algebra equipped with an invertible module with involution $M$ such that the involution on $M$ is $R$-sesquilinear then $(\Dperf(B),\QF^{\gs}_M)$ is a module over $(\Dperf(B),\QF^{\gs}_{\sig})$, and if $\lambda$ is any form parameter over $B$ refining $M$ then $(\Dperf(B),\QF^{\g\lambda}_M)$ is a module over $(\Dperf(R),\QF^{\g\fr}_{\sig})$. In general, $(\Dperf(B),\QF^{\g\lambda}_M)$ is not a module over $(\Dperf(R),\QF^{\gs}_{\sig})$, for instance $(\Dperf(\ZZ),\QF^{\g\fr}_\ZZ)$ is not a module over $(\Dperf(\ZZ),\QF^\gs_\ZZ)$, see also Remark~\reffour{remark:criterion-genuine-symmetric}.

In particular, $\GW(\Dperf(B),\QF^{\g\lambda}_M)$ and $\L(\Dperf(B),\QF^{\g\lambda}_M)$ are modules over $\GW(\Dperf(R),\QF^{\g\fr}_\sig)$ and $\L(\Dperf(R),\QF^{\g\fr}_\sig)$, but in general not over $\GW(\Dperf(R),\QF^{\gs}_\sig)$ and $\L(\Dperf(R),\QF^{\gs}_\sig)$. For instance, $\GW^{\g\fr}(\ZZ)$ is not a module over $\GW^{\gs}(\ZZ)$, and similarly $\L^{\g\fr}(\ZZ)$ is not a module over $\L^{\gs}(\ZZ)$; indeed it it were a module, since $\L^{\g\fr}_0(\ZZ) \to \L^{\gs}_0(\ZZ)\cong \ZZ$ is surjective, we would obtain a splitting $\L^{\g\fr}(\ZZ) \simeq \L^\gs(\ZZ) \oplus M$ contradicting that $\L^{\g\fr}_1(\ZZ) = 0$ and $\L^{\gs}_1(\ZZ) = \ZZ/2$; a similar argument also shows the case of $\GW$.
\end{example}



Next, we aim to describe more explicitly the multiplicative structure on some GW- and L-spectra, focussing first on ring structures on homotopy groups.
\begin{example}
\label{example:universal}%
The most fundamental example of a symmetric monoidal Poincaré $\infty$-category is the unit $(\Spaf,\QF^{\uni})$ of $\Catp$ and we obtain $\Einf$-ring spectra $\GW^\uni(\SS)$ and $\L^\uni(\SS)$ over which all GW- and L-spectra are modules. By Proposition~\reftwo{proposition:pi0gwuniv}, we have
$\GW_0^\uni(\SS) = \ZZ[e,h]/(e^2-8e,he-8h,h^2-2h)$ and $\L_0^\uni(\SS) = \ZZ[e]/(e^2-8e)$.
\end{example}

\begin{example}
\label{example:L-ring-Z}%
Another fundamental commutative ring is $\ZZ$. Let us record here the multiplication on various forms of its L-groups, due to Mishchenko, Kervaire--Milnor, and Ranicki, see \cite{Ranickiyellowbook}*{Section 4.3.1} or Corollary~\refthree{corollary:compute-L-groups-global-field} and Proposition~\refthree{proposition:compute-quadratic-L-groups}. Namely, we have $\L^\s_*(\ZZ) = \ZZ[x^{\pm 1},e]/(2e,e^2)$ with $|x|=4$ and $|e|=1$ and 
\[
\L^\qdr_n(\ZZ) = \begin{cases} 8\ZZ & \text{ for } n\equiv 0 (4) \\ 0 & \text{ for } n \equiv 1,3 (4) \\ \ZZ/2 & \text{ for } n \equiv 2(4) \end{cases}.
\]
which is an $\L_*^\s(\ZZ)$-module in a unique fashion (the module multiplication by $e$ is trivial for degree reasons). To determine the genuine symmetric L-ring of $\ZZ$, recall from \paperthree that the ring map $\L^\gs_*(\ZZ) \to \L^\s_*(\ZZ)$ is an isomorphism for $*\geq -2$ and the $\L^\gs_*(\ZZ)$-module map $\L^\qdr_*(\ZZ) \to \L^\gs_*(\ZZ)$ is an isomorphism for $*\leq -3$. Since the $\L^\gs(\ZZ)$-module structure on $\L^\qdr(\ZZ)$ is restricted from its canonical $\L^\s(\ZZ)$-module structure, one can infer from the above the ring structure on $\L^\gs_*(\ZZ)$, see \cite{HLN}*{Remark 6.3} for explicit formulae.
\end{example}

In fact, the computation of the symmetric and quadratic L-groups of $\ZZ$ together with the canonical module structure explained above can be generalised as follows.

\begin{example}
\label{example:L-rings-dedekind}%
Let $D$ be a Dedekind domain whose fraction field $F$ is a global field of characteristic different from 2, e.g.\ that $D$ is a number ring. We have computed in Corollary~\refthree{corollary:compute-L-groups-global-field} that its symmetric L-groups are 4-periodic and given by
\[
\L^\s_n(D) = \begin{cases} \W^\s(D) & \text{ for } n \equiv 0 (4),\\
						\bigoplus_{\I} \ZZ/2 & \text{ for } n \equiv 1 (4),\\
						0 & \text{ for } n \equiv 2 (4), \\
						\Pic(D)/2 & \text{ for } n \equiv 3 (4), \end{cases}
\]
where $\I$ is the set of dyadic primes in $D$. In fact, for $\mathfrak{p} \in \I$ such a prime, the map $D \to D_\mathfrak{p}$ induces the associated projection on $\L_1^\s(-)$. There are now three types of basic products in the ring $\L_\ast^\s(D)$, described by $\W^\s(D)$-module structures on $\Pic(D)/2$ and $\bigoplus_\I \ZZ/2$, respectively, and the product
\[
\Pic(D)/2 \times \bigoplus_\I \ZZ/2 \to \W^\s(D).
\]
We claim that this product vanishes. To show this, we may postcompose with the injection $\W^s(D) \to \W^\s(F)$, so that the claim follows from the naturality of the product and the fact that $\Pic(F) = 0$. Next, we claim that the two $\W^\s(D)$-module structures on $\L_{1}^\s(D)$ and $\L_{-1}^\s(D)$ are obtained from their canonical $\FF_2$-module structures via the rank mod 2 homomorphism $\W^\s(D) \to \FF_2$. We first argue this in case of $\L_1^\s(D)$. From the above concrete description of this group, we may assume that $D=D_\mathfrak{p}$ is local for some dyadic prime $\mathfrak{p}$. Since the endomorphism ring of $\ZZ/2$ is $\FF_2$, the claimed module structure is equivalently described by a ring homomorphism $\W^\s(D_\mathfrak{p}) \to \FF_2$, of which the rank mod 2 homomorphism is the unique one. Indeed, $\W^\s(D_{\mathfrak{p}})$ is generated by units represented by the 1-dimensional forms $\langle a \rangle$ for $a \in D_{\mathfrak{p}}^\times$, see \cite{knebusch-rosenberg-ware2}*{Theorem 1.16}, so it follows that $\langle a \rangle$ is sent to $1$ under any ring homomorphism $\W^\s(D_{\mathfrak{p}}) \to \FF_2$. It follows that any such map is the rank mod 2 homomorphism as claimed.
Finally, we need to argue that the $\W^\s(D)$-module structure on $\Pic(D)/2$ is also obtained via the rank mod 2 homomorphism $\W^\s(D) \to \FF_2$. In Example~\reffour{example:free-dedekind} we have shown that the composite of maps of $\Einf$-ring spectra
\[
\L^\s(D) \to \K(D)^{\tC} \to \K_0(D)^{\tC}
\]
induces an isomorphism on $\pi_{-1}$. The result then follows from the well-known computation that $\K_0(D)$ is isomorphic to the trivial square zero extension $\ZZ \oplus \Pic(D)$. Indeed, recall that 
\[
\widehat{\mathrm{H}}^n(\Ct;\K_0(D)) = \begin{cases} \FF_2 \oplus \Pic(D)[2] & \text{ for } n \equiv 0 (2) \\ \Pic(D)/2 & \text{ for } n \equiv 1(2) \end{cases}
\]
so that the lax symmetric monoidality of Tate cohomology implies that the product
\[
\Pic(D)[2] \times \Pic(D)/2 \subseteq \widehat{\mathrm{H}}^0(\Ct;\K_0(D)) \times \Pic(D)/2 \to \Pic(D)/2
\]
vanishes as claimed.

Next, recall that the quadratic L-groups of $D$ are given as follows
\[
\L^\qdr_n(D) = \begin{cases} \W^\qdr(D) & \text{ for } n \equiv 0 (4),\\
						0 & \text{ for } n \equiv 1 (4), \\
						\bigoplus_{\I} \W^\qdr(\FF_\mathfrak{p}) & \text{ for } n \equiv 2 (4),\\
						A \rtimes \Pic(D)/2 & \text{ for } n \equiv 3 (4) \end{cases}
\]
where $\FF_{\mathfrak{p}} = D/\mathfrak{p}$ is the residue field at $\mathfrak{p}$, $A$ is the cokernel of the map $\W^\s(D) \oplus \W^\qdr(D^\cwedge_2) \to \W^\s(D^\cwedge_2)$ and $A \rtimes \Pic(D)/2$ is a certain extension of $\Pic(D)/2$ by $A$, see Proposition~\refthree{proposition:compute-quadratic-L-groups}. Moreover, the projection $\L_2^\qdr(D) \to \W^\qdr(\FF_{\mathfrak{p}})$ is induced by the projection $D \to \FF_\mathfrak{p}$.
Let us assume that $2$ is not invertible in $D$ (and hence the fraction field of $D$ is a number field), as else the map $\L^\qdr(D) \to \L^\s(D)$ is an equivalence and we have already computed all relevant products. 

The $\L^\s_*(D)$-module structure on $\L^\qdr_*(D)$ is now determined by various multiplications. Using similar arguments as above, one can show that these products are given as follows:

The $\W^\s(D)$-module structures on $\L_n^\qdr(D)$ are the evident ones for $n=0,1$, and is the restricted $\FF_2$-module structure on $\L^\qdr_2(D)$ obtained from the rank mod 2 homomorphism $\W^\s(D) \to \FF_2$. Moreover, the extension $A \rtimes \Pic(D)/2$ is one of $\W^\s(D)$-modules, where $\Pic(D)/2 = \L_{-1}^\s(D)$ is again a $\W^\s(D)$-module via the rank mod 2 map $\W^\s(D) \to \FF_2$ as above, and where the surjection $\L_0^\s(D^\cwedge_2) \to A$ is also one of $\W^\s(D)$-modules.
Next, we describe the products with $\L_1^\s(D)$ and $\L_{-1}^\s(D)=\Pic(D)/2$. We begin with $\L_1^\s(D)$ in which case there are products
\[
\L_1^\s(D) \times \L^\qdr_2(D) \to \L^\qdr_3(D) \quad \text{ and } \quad \L_1^\s(D) \times \L^\qdr_3(D) \to \W^\qdr(D).
\]
of which the latter vanishes (for this, one may use that the symmetrisation map $\W^\qdr(D) \to \W^\s(D)$ is injective, see the proof of Corollary~\refthree{corollary:finite-generation} and the previously established products on $\L^\s_*(D)$). 
For the first product, the composite with the surjection $\L^\qdr_3(D) \to \L^\s_3(D)=\Pic(D)/2$ vanishes, so it is determined by a map
\[
\L_1^\s(D) \times \L_2^\qdr(D) \to A
\]
about which we do not know more at this point.
Similarly, there are products
\[
\Pic(D)/2 \times \W^\qdr(D) \to \L_{3}^\qdr(D) \quad \text{ and } \quad \Pic(D)/2 \times \L^\qdr_3(D) \to \L^\qdr_2(D).
\]
of which the latter vanishes again. 
Moreover, the former product followed by the surjection $\L^\qdr_3(D) \to \L^\s_3(D)=\Pic(D)/2$ vanishes, and is hence determined by a product 
\[
\Pic(D)/2 \times \W^\qdr(D) \to A
\]
about which we again do not know more at this point.
Indeed, the composite under investigation is restricted form the tautological $\FF_2$-module structure of $\Pic(D)/2$ via the composite $\W^\qdr(D) \to \W^\s(D) \to \FF_2$ where the latter map is the rank mod 2 homomorphism. Now since we have assumed that $2$ is not invertible in $D$ there exists at least one dyadic prime $\mathfrak{p}$ and the  composite factors through the map $\W^\qdr(\FF_\mathfrak{p}) \to \W^\s(\FF_\mathfrak{p})$ which vanishes, see \cite{milnor-symmetric}*{Appendix 1}. 

Finally one can use the comparison results of \paperthree and \cite{Schlichting-genuine}, to find that the map $\L^\gs_*(D) \to \L^\qdr_*(D)$ is an isomorphism for $* \geq -2$ and the map $\L^\qdr_*(D) \to \L^\gs_*(D)$ is one for $* \leq -3$, so the above also determines (up to the ambiguity involving products with image in $A$) the ring structure on $\L^\gs_*(D)$.
\end{example}






Next, we will describe several GW- and L-spectra as $\Einf$-rings in terms of more familiar $\Einf$-rings. The first two examples are of the kind where the comparison maps to homotopy or Tate fixed points on K-theory are in fact (close to) equivalences:

\begin{example}
\label{example:perfect-fields-of-characteristic-2}%
For any perfect field of characteristic 2, the map of $\Einf$-ring spectra $\L^\s(F) \to \K(F)^{\tC}$ is an equivalence, see Remark~\refthree{remark:K-of-perfect-field}. Moreover, part of the argument in loc.\ cit.\ is that the map $\K(F) \to \ZZ$ is a 2-adic equivalence and hence induces an equivalence of $\Einf$-ring spectra $\K(F)^{\tC} \to \ZZ^{\tC}$. As a result, we obtain an $\Einf$-equivalence $\L^\s(F) \simeq \ZZ^{\tC}$, and deduce that the map of $\Einf$-ring spectra $\GW^\s(F) \to \ZZ^{\hC}$ is an equivalence after 2-adic completion. We also note that there exists a canonical equivalence of $\Einf$-ring spectra $\mathrm{THH}(\FF_2) \simeq \tau_{\geq0}\ZZ^{\tC}$ whose source is canonically an $\Einf$-$\FF_2$-algebra, see \cite{NS}*{Corollary IV.4.16}. In particular, we deduce that $\L^\s(A)$ is acquires to structure of an $\Einf$-$\FF_2$-algebra whenever $A$ is itself an $\Einf$-$\FF_2$-algebra.
\end{example}

\begin{example}
\label{ex:algebraically-closed-fields}%
For a quadratically closed field $K$ of characteristic different from 2, we argue that $\L^\s(K)$ is canonically equivalent to $\ku^{\tC}$, where the $\Ct$-action on $\ku$ is via complex conjugation. To see this, first note that $\L^\s_*(K) = \FF_2[u^{\pm1}]$ with $|u|=4$. Indeed, $\L^\s_*(K)$ is 4-periodic and vanishes in degrees not divisible by $4$, while $\L_0^\s(K) \cong \W_0^\s(K) \cong \FF_2$, see Remark~\refthree{remark:odd-vanishing}
and~\cite{milnor-symmetric}*{\S 3 Remark 3.4}. If $K$ has characteristic $0$ then there is an algebraically closed field $L$ which contains $K$ and $\CC$ as subfields. The maps $\L^\s(K) \to \L^\s(L) \leftarrow \L^\s(\CC)$ are then equivalences of $\Einf$-rings, because all homotopy rings are isomorphic to the graded ring $\FF_2[u^{\pm1}]$ with $|u|=4$, which has only the identity as endomorphism. If $K$ has positive characteristic different from $2$ we may choose an algebraic closure $\overline{K}$ of $K$ and the map $\L^\s(K) \to \L^\s(\overline{K})$ is an equivalence of $\Einf$-rings for the same reason as before. Now $\overline{K}$ is a perfect field of characteristic $p >2$, so we may consider the complete discrete valuation ring $D$ of $p$-typical Witt vectors of $\overline{K}$ and denote by $F$ an algebraic closure of the fraction field of $D$, an algebraically closed field of characteristic $0$. We claim that both maps $\L^\s(\overline{K}) \leftarrow \L^\s(D) \to \L^\s(F)$ are equivalences of $\Einf$-rings. To see that the left map is an equivalence, we note that $2$ is invertible in $D$ and $\overline{K}$, so that the map is equivalent to the map $\L^\qdr(D) \to \L^\qdr(\overline{K})$ which is an equivalence by Proposition~\refthree{proposition:hensel} since $D$ is $p$-adically complete and $D/p \to K$ is an isomorphism. In particular, $\L^\s_*(D)$ is also isomorphic to $\FF_2[u^{\pm1}]$ with $|u|=4$. The same argument as before then implies that also the map $\L^\s(D) \to \L^\s(F)$ is an equivalence. We have therefore now argued that any two quadratically closed fields of characteristic different from $2$ have canonically equivalent symmetric L-spectra. Note then that there are canonical maps of $\Einf$-rings $\L^\s(\CC) \to \K(\CC)^{\tC} \to \ku^{\tC}$. By \cite{Greenlees-reality}*{Corollary 3.3}, $\pi_*(\ku^{\tC})$ is again isomorphic to $\FF_2[u^{\pm1}]$ with $|u|=4$, so our favourite argument implies that the map $\L^\s(\CC) \to \ku^{\tC}$ is an equivalence of $\Einf$-rings.

As for $\GW$-theory, if $K$ is separably closed of characteristic different from 2 then it follows from Suslin's rigidity results that there is an 
equivalence of $\Einf$-rings $\K(K)^\cwedge_2 \simeq \ku^\cwedge_2$, and the above implies that $\GW^\s(K)^\cwedge_2 \to (\ku^\cwedge_2)^{\hC}$ is an equivalence. In particular, we obtain an equivalence of $\Einf$-rings $\tau_{\geq0} \GW^\s(K)^\cwedge_2 \simeq \ko^\cwedge_2$ for $K$ separably closed of characteristic different from $2$.
\end{example}

To move on, let us recall that for a commutative ring with involution $(R,\sigma)$ and $\sigma$-fixed unit $u \in R^\times$, there is associated the 1-dimensional $\sigma$-symmetric form $(x,y) \mapsto xu\sigma(y)$, giving rise to a unit $\langle u \rangle$ in the classical symmetric Grothendieck-Witt ring $\GW_0^\s(R,\sigma)$ of $(R,\sigma)$. Since multiplication by $u$ provides an isomorphism between the forms associated to $1$ and to $u\sigma(u)$, this construction gives rise to a map of commutative rings 
\[
\ZZ\Bigl[\frac{(R^\times)^\sigma}{\mathrm{N}(R^\times)}\Bigr] \to \GW_0^\s(R,\sigma),
\]
where $\mathrm{N}(R^\times)$
refers to the subgroup given by elements of the form $u\sigma(u)$. This map is instrumental to many of the classical descriptions of the symmetric Grothendieck-Witt and Witt rings of commutative rings. The following provides a lift of this algebraic construction to the level of the $\Einf$-ring spectrum $\GW^\s(R,\sigma)$. We will then use this to lift several known algebraic descriptions of 
symmetric Grothendieck-Witt and Witt rings
to a description of the $\Einf$-ring structure on suitable symmetric Grothendieck--Witt and L-spectra.

\begin{construction}
\label{construction:Bgl-to-gl}%
Let $(\C,\QF)$ be a symmetric monoidal Poincar\'e $\infty$-category. Since the map $\Poinc \to \GWspace$ is a lax symmetric monoidal transformation, we find that the map $\Poinc(\C,\QF) \to \GWspace(\C,\QF)$ is an $\Einf$-map with respect to tensor product on each side. Passing to $\otimes$-invertible components, we obtain a natural map of $\Einf$-groups $\Poinc(\C,\QF)^\times \to \gl_1(\GW(\C,\QF))$.
In the case where $(\C,\QF) = (\Modp{R},\QF^{\sym}_{\sig})$ for an $\Einf$-ring spectrum with involution $(R,\sig)$ %
this gives a map $\Pic(R)^{\hC} \to \gl_1(\GW^{\sym}(R,\sig))$,
which fits into a commutative diagram
\[
\begin{tikzcd}
	\Pic(R)^{\hC} \ar[r] \ar[d] & \gl_1(\GW^\s(R,\sig)) \ar[r] & \gl_1(\K(R,\sig)^{\hC}) \ar[d] \\
	\Pic(R) \ar[rr] && \gl_1(\K(R))
\end{tikzcd}
\]
whose lower horizontal map is well-known in algebraic K-theory.
Now suppose that $R$ is discrete. The inclusion of the identity component $\B R^{\times} = \mathrm{Bgl}_1(R) \subseteq \Pic(R)$
then induces a map
\[
(\B R^\times)^{\hC} \to \Pic(R,\sig)^{\hC} \to \gl_1(\GW^\s(R,\sig))
\]
whose first map is an equivalence if the canonical map $\ZZ \to \pi_0(\Pic(R))$ is an isomorphism.
Observe now that $\B R^\times$ is the space underlying the $\ZZ$-module spectrum $\Sigma \GEM R^\times$ and that the $\Ct$-action induced by $\sigma$ on it is automatically $\ZZ$-linear; in fact, it is induced by the $\Ct$-action $u \mapsto \sigma(u)^{-1}$ on the abelian group $R^\times$. It follows that also the fixed point spectrum $(\B R^\times)^{\hC}$ is canonically a $\ZZ$-module, and hence (non-canonically) splits as $\Einf$-group into a direct sum $\mathrm{H}^1(\Ct;R^\times)\oplus \B \mathrm{H}^0(\Ct;R^\times)$.
The choice of such a splitting then induces an $\Einf$-map 
\[
\frac{(R^\times)^\sigma}{\mathrm{N}(R^\times)} = \mathrm{H}^1(\Ct;R^\times) \to \Pic(R)^{\hC} \to \gl_1(\GW^\s(R,\sig)),
\]
and by adjunction, %
maps of $\Einf$-ring spectra
\[
\SS\Bigl[\frac{(R^\times)^\sigma}{\mathrm{N}(R^\times)}\Bigr] \to \GW^\s(R,\sig) \to \L^\s(R,\sig).
\]
On $\pi_0$, this map sends the class of a $\sig$-invariant unit $u \in R^\times$ to the 1-dimensional symmetric form $\langle u \rangle$. In particular, if $R$ is equipped with the identity involution, there exist maps of $\Einf$-ring spectra
\[
\SS\Bigl[\frac{R^\times}{(R^\times)^2}\Bigr] \to \GW^\s(R) \to \L^\s(R)
\]
which on $\pi_0$ send a unit $u \in R^\times$ to $\langle u \rangle$.
\end{construction}


\begin{example}
In \paperthree, we have utilised a topological version of Grothendieck-Witt theory $\GW_\topo ^\s(\RR)$ introduced in the context of Banach algebras with involution by Schlichting in \cite{schlichting-derived}*{\S 10} and which also features in the analysis of symmetric L-theory of real $C^*$-algebras in \cite{LNS}. By construction, it comes with $\Einf$-ring map $\GW^\s(\RR) \to \GW^\s_\topo(\RR)$ and by \cite{LNS}*{Remark 4.2} it also comes with an $\Einf$-ring map $\ko \to \GW^\s_\topo(\RR)$; informally this map is induced by sending a finite dimensional real vector space to its (essentially unique) positive definite symmetric bilinear form. As a consequence of this, Construction~\reffour{construction:Bgl-to-gl}, and the observation that $\RR^\times/(\RR^\times)^2 = \{\pm 1\} = \Ct$, there exists a map of $\Einf$-rings
\[
\ko[\Ct] \to \tau_{\geq0}\GW_\topo^\s(\RR)
\]
which, by the argument used in the proof of \cite{LNS}*{Corollary 3.6}, is an equivalence.
\end{example}

\begin{example}
Let $K$ be a complete discretely valued field with ring of integers $D$, a complete discrete valuation ring with residue field $k$. Assume that the characteristic of $k$ is different from $2$ and choose a uniformiser $\pi \in D$. 
Consider the map $\Ct \to K^\times/(K^\times)^2$ with image the subgroup generated by $\pi \in D$. Applying Construction~\reffour{construction:Bgl-to-gl}, we deduce the existence of a map of $\Einf$-rings $\L^\s(D)[\Ct] \to \L^\s(K)$ which on $\pi_0$ sends the generator of $\Ct$ to $\langle \pi \rangle$.
As in Example~\reffour{ex:algebraically-closed-fields}, the projection $D \to k$ induces an equivalence $\L^\s(D) \to \L^\s(k)$ since $D$ is $\pi$-adically complete and $D/\pi \to k$ is an isomorphism: Indeed, since $2$ is invertible in $D$ and $k$, it suffices to show that $\L^\qdr(D) \to \L^\qdr(k)$ is an equivalence, which follows from the $\pi$-completeness of $D$, see e.g.\ Proposition~\refthree{proposition:hensel}. In total, we obtain a map of $\Einf$-ring spectra
\[
\L^\s(k)[\Ct] \to \L^\s(K)
\]
which induces the Springer decomposition $\W^\s(k)[\Ct] \cong \W^\s(K)$ of symmetric Witt groups on $\pi_0$, see e.g.\ \cite{Lam}*{VI Corollary 1.7}. Since both homotopy rings are 4-periodic, with non-trivial homotopy only in degrees divisible by $4$, see again Remark~\refthree{remark:odd-vanishing}, the above map is an equivalence of $\Einf$-rings.
For instance, we obtain that for $p$ an odd prime, there is an equivalence of $\Einf$-ring spectra 
\[
\L^\s(\FF_p)[\Ct] \simeq \L^\s(\QQ_p).
\]
That the same result is not true for $p=2$ is well-known and already fails on the level of $\pi_0$; Here the source has 4 elements while the target has $32$-elements.
\end{example}

\begin{example}
\label{example:Gm}%
Let $R$ be a regular noetherian ring of finite Krull dimension. Consider the ring of Laurent polynomials $R[t^{\pm1}]$ as a commutative ring with trivial involution as in Remark~\reffour{remark:all-involutions}. Similarly as in the example above, consider the map 
\[
\Ct \to \Ct \times \Ct \cong \ZZ[t^{\pm1}]^\times/(\ZZ[t^{\pm1}]^\times)^2 \to R[t^{\pm1}]^\times/(R[t^{\pm1}]^\times)^2
\]
with image the subgroup generated by $t$. From Construction~\reffour{construction:Bgl-to-gl}, we obtain maps of $\Einf$-rings 
\[
\GW^\s(R)[\Ct] \to \GW^\s(R[t^{\pm1}]) \quad \text{and} \quad \L^\s(R)[\Ct] \to \L^\s(R[t^{\pm1}]).
\]
Under the the equivalences indicated in Remark~\reffour{remark:all-involutions}, one can show that the underlying map of spectra for the latter map is the identity while for $\GW$ it is induced by the boundary map in the Genauer sequence (which is not an equivalence unless $R = 0$).


In \S\reffour{subsection:mult-shaneson} below, we will discuss the multiplicative properties of the Ranicki-Shaneson decompositions for the L-theory of the group ring $R[\ZZ]$, with involution induced by the group inversion on $\ZZ$.
\end{example}

We deduce the following computation from Example~\reffour{example:Gm}.
\begin{example}
Let $R$ be a regular noetherian commutative ring of finite Krull dimension and denote by $\mathbb{P}^1(R)$ the projective line over $R$. In \cite{motives} it is shown that the square
\[
\begin{tikzcd}
	\Perf(\mathbb{P}^1(R)) \ar[r]\ar[d] & \Perf(\mathbb{A}^1(R)) \ar[d] \\
	\Perf(\mathbb{A}^1(R)) \ar[r] & \Perf(\mathbb{G}_m(R))
\end{tikzcd}
\]
canonically refines to a Poincar\'e-Verdier square using the symmetric Poincar\'e structures throughout. Moreover, it is shown that $\GW^\s$ and $\L^\s$ are $\mathbb{A}^1$-invariant in this situation, so that there are pullback squares
\[
\begin{tikzcd}
	\GW^\s(\mathbb{P}^1(R)) \ar[r] \ar[d] & \GW^\s(R) \ar[d] && \L^\s(\mathbb{P}^1(R)) \ar[r] \ar[d] & \L^\s(R) \ar[d] \\
	\GW^\s(R) \ar[r] & \GW^\s(\mathbb{G}_{m}(R)) && \L^\s(R) \ar[r] & \L^\s(\mathbb{G}_{m}(R))
\end{tikzcd}
\]
of $\Einf$-ring spectra. Through the canonical section $1 \colon R \to \mathbb{G}_m(R) \subseteq \mathbb{P}^1(R)$, these diagrams are in fact pullback diagrams of augmented $\Einf$-$\GW^\s(R)$- and $\L^\s(R)$-algebras, respectively. Example~\reffour{example:Gm} implies that $\L^\s(\mathbb{P}^1(R))$ as an $\Einf$-ring spectrum is given by $\L^\s(R) \otimes P$ where $P = \SS \times_{\SS[\Ct]} \SS$ is the loop object of $\SS[\Ct]$. If $R$ is a commutative $\FF_2$-algebra, recall from Example~\reffour{example:perfect-fields-of-characteristic-2} that $\L^\s(R)$ admits a canonical $\Einf$-$\FF_2$-algebra structure. In particular, as an $\Einf$-ring, we have $\L^\s(\mathbb{P}^1(R)) \simeq \L^\s(R) \otimes_{\FF_2}[ \FF_2 \times_{\FF_2[\Ct]} \FF_2]$ and $\FF_2[\Ct]$ is a trivial square-zero extension. Consequently, if $R$ is a commutative $\FF_2$-algebra, we deduce that $\L^\s(\mathbb{P}^1(R))$ is equivalent to the trivial square-zero extension $\L^\s(R) \oplus \Omega\L^\s(R)$ as $\Einf$-$\L^\s(R)$-algebra.

Moreover, by \cite{formality}*{Theorem B}, it follows that as $\Eone$-$\GW^\s(R)$-algebra, $\GW^\s(\mathbb{P}^1(R))$ is the trivial square-zero extension on $\GW(R,\Omega \QF^\s)$. Similarly, as $\Eone$-$\L^\s(S)$-algebra, $\L^\s(\mathbb{P}^1(R))$ is the trivial square-zero extension on $\Omega \L^\s(R)$, also when $R$ is not an $\FF_2$-algebra.


As a concrete example, consider the projective line $\mathbb{P}^1(D)$ over the ring of integers in a number field. Then we obtain an equivalence of augmented $\Eone$-rings $\L^\s(\mathbb{P}^1(D)) \simeq \L^\s(D) \oplus \Omega \L^\s(D)$, where the latter is the trivial square-zero extension. In particular, by Corollary~\refthree{corollary:compute-L-groups-global-field}, we find 
\[
\L^{\sym}_0(\mathbb{P}^1(D)) = \L^{\sym}_0(D) \oplus \L^{\sym}_1(D) = \W^\s(D) \oplus (\ZZ/2)^{\oplus n}
\]
where $n$ is the number of dyadic primes in $D$ and the multiplication on $(\ZZ/2)^n$ is trivial, and the $\W^\s(D)$-module structure on $(\ZZ/2)^{\oplus n}$ is via the rank mod 2 homomorphism $\W^\s(D) \to \FF_2$ as explained in Example~\reffour{example:L-rings-dedekind}.
\end{example}




\begin{remark}
As a consequence of the $\Einf$-ring map $\SS[R^\times/(R^\times)^2] \to \GW^\s(R)$ from Construction~\reffour{construction:Bgl-to-gl}, all the units $\langle u \rangle$ in $\GW_0^\s(R)$ are \emph{strict} units in the $\Einf$-ring $\GW^\s(R)$, i.e.\ are represented by $\Einf$-maps $\SS[\ZZ] \to \GW^\s(R)$, and the same therefore holds true for their images in $\L_0^\s(R)$. In fact, these strict units are strictly of order $2$, that is, their classifying maps factor as the composite of $\Einf$-ring maps $\SS[\ZZ] \to \SS[\Ct] \to \GW^\s(R) \to \L^\s(R)$. 

Strict units allow for more constructions than general units. For instance, for an $\Einf$-ring spectrum $A$ and $u \in \pi_0(A)$ a strict unit classified by an $\Einf$-map $\SS[\ZZ] \to A$, one may consider the base-change $A \otimes_{\SS[\ZZ]} \SS$ of $A$ along the augmentation $\SS[\ZZ] \to \SS$ induced by the group homomorphism $\ZZ \to \{1\}$ on spherical group rings. One can compute that the underlying $A$-module of this base-change is given by $A/(u-1)$, so that this $A$-module is canonically an $\Einf$-$A$-algebra. In fact, more generally, one can show that $A/(u-1)^k$ is an $\Einf$-$A$-algebra for all $k\geq 1$: Considering the universal case $\SS[\ZZ]$ with strict unit $t$ corresponding to the generator $1 \in \ZZ$, we find that the fibre of the augmentation map $\SS[\ZZ] \to \SS$ is a Smith ideal in $\SS[\ZZ]$. Taking (pointwise) $k$th powers of the Smith ideal gives a new Smith ideal, whose quotient can be shown to deserve the name $\SS[t]/(t-1)^k$. Then $A/(u-1)^k$ is given by $A \otimes_{\SS[\ZZ]} \SS[t]/(t-1)^k$.

Now, it follows from the above construction that for a discrete commutative ring $R$, the element $\langle -1 \rangle = -1 \in \L_0^\s(R)$ is a strict unit. In particular, for all $k\geq 1$, we deduce that $\L^\s(R)/2^k$ is canonically an $\Einf$-ring spectrum, in fact an $\Einf$-algebra over $\L^\s(\ZZ)/2^k$. Moreover, one can form the base-change $\L^\s(R) \otimes_{\SS[\Ct]} \SS$ of $\L^\s(R)$ along the augmentation of $\SS[\Ct]$ which is again an $\Einf$-$\L^\s(R)$-algebra whose underlying $\L^\s(R)$-module can be shown to be given by $\L^\s(R)_{\hC}$, where the $\Ct$-action is induced by multiplication by the strict unit of order 2 given by $-1$.


As a consequence of the $\Einf$-equivalences $\L^{\sym}(\ZZ)[1/2] = \L^{\sym}(\RR)[1/2] = \KO[1/2]$, the second arising from \cite{LN}*{Theorem 5.2}, $-1 \in \pi_0(\KO[1/2])$ is also given the structure of a strict unit. This is not surprising, though: Indeed, it follows from results of Carmeli that the spectrum of strict units in $\SS[1/2]$ is discrete with $\pi_0$ given by $\{\pm 1\}$, while the spectrum of strict units of $\SS$ is contractible \cite{Carmeli}*{Theorem B}. In particular, in general, $-1 \in \pi_0(A)$ does not admit the structure of a strict unit.
\end{remark}




















 








\subsection{The periodicity elements of Karoubi and Ranicki}
\label{subsection:PeriodicityviaMultiplication}%


In this section we aim to explain the relation between the multiplicative structure on Grothendieck-Witt- and $\L$-spectra and their periodicity properties, and then use this to reconcile our formulations of these in \paperone and \papertwo with Ranicki's formulation of his periodicity isomorphisms for $\L$-groups and Karoubi's original formulation of his fundamental theorem. In particular, we will show that the equivalences
\[
\mathbb S^2 \otimes \L^\sym(R) \simeq \L^{-\sym}(R) \quad \text{and} \quad \mathbb S^1 \otimes \mathrm V^\sym(R) \simeq \mathrm U^{-\sym}(R),
\]
produced in \papertwo are induced by multiplication with certain periodicity classes $\rho$ and $\kappa$ along with appropriate generalisations to form parameters and $\Einf$-rings, where %
$\mathrm V$ and $\mathrm U$ are given by
\[
\mathrm V(\C,\QF) = \fib(\GW(\C,\QF) \xrightarrow{\fgt} \K(\C))  \quad \text{and} \quad \mathrm U(\C,\QF) = \fib(\K(\C) \xrightarrow{\hyp} \GW(\C,\QF)),
\]
see \S\reftwo{subsection:Genauer-Bott-Karoubi}.

To get started, %
recall that by Corollary~\refone{corollary:quadratic-genuine} any hermitian structure $\QF \colon \C\op \rightarrow \Sp$ canonically lifts to a functor $\tilde\QF\colon \C\op \to \Sp^{\g\Ct}$ to genuine $\Ct$-spectra along the genuine fixed point functor $(-)^\gC \colon \Sp^\gC \rightarrow \Sp$. The underlying $\Ct$-spectrum of $\tilde \QF(X)$ is given by $\Bil_\QF(X,X)$ and its geometric fixed points by $\Lambda_\QF(X)$. We can now tensor this lift valuewise with the genuine representation sphere $\SS^{m+n\sigma} \in \Sp^{\g\Ct}$ for $m,n \in \mathbb Z$ and apply genuine fixed points. Slightly abusing notation we will denote the resulting functor by $\QF\qshift{m+n\sigma}$. Inspection of the isotropy separation square of $\mathbb S^{[m+n\sigma]} \otimes \tilde\QF$ then immediately verifies that $\QF\qshift{m+n\sigma}$ is again a hermitian structure on $\C$ with
\[
\Bil_{\QF\qshift{m+n\sigma}} \simeq \SS^{m+n\sigma} \otimes \Bil_\QF \quad \text{and} \quad \Lambda_{\QF\qshift{m+n\sigma}} \simeq \mathbb S^m \otimes \Lambda_\QF,
\]
connected by the composite
\[
\SS^m \otimes \Lambda_\QF \longrightarrow \SS^m \otimes (\Bil_\QF \circ \Delta)^\tC \simeq (\SS^{m+n\sigma} \otimes \Bil_\QF \circ \Delta)^\tC,
\]
where the first map is the image of the glueing map for $\QF$ and the second induced by the inclusion $\SS^0 \rightarrow \SS^{n\sigma}$ as fixed points. In particular, the operation $(-)\qshift{m+n\sigma}$ preserves both symmetric and quadratic Poincaré structures.

To obtain the interaction of this construction with the tensor product of Poincaré $\infty$-categories, we note the following alternative description. Theorem~\refone{theorem:classification-genuine-modules} and Remark~\refone{remark:NormModules} imply that $\Funq(\Spaf) \simeq \Sp^{\g\Ct}$ as symmetric monoidal $\infty$-categories, see \S\refone{subsection:examples-monoidal}). In particular, we may identify hermitian structures on $\Spaf$ with genuine $\Ct$-spectra by regarding their classifying squares as %
isotropy separation ones. Consequently, the symmetric monoidal structure on $\Cath$ determines an action of $\Sp^{\g\Ct}$ on $\Funq(\C)$ for any stable $\infty$-category $\C$. Unwinding definitions one finds that acting with $\mathbb S\qshift{m+n\sigma}$ on some $\QF$ reproduces exactly $\QF\qshift{m+n\sigma}$. %
We therefore have that
\[
(\C,\QF\qshift{m+n\sigma}) \otimes (\D,\Psi\qshift{k+l\sigma}) = (\C \otimes \D,(\QF \otimes \Psi)^{(m+k)+(n+l)\sigma}).
\]
In particular, the operation $(-)\qshift{m+n\sigma}$ is invertible with inverse $(-)\qshift{-m-n\sigma}$.  %
Now for the periodicity elements:

\begin{construction}
We will first define two elements 
\[
\kappa \in \Omega^\infty\GW(\Spaf,(\QF^\uni)\qshift{-1-\sigma}) \quad \text{and} \quad \delta \in \Omega^{\infty-1}\GW(\Spa^\fin,(\QF^\uni)\qshift{-1}),
\]
multiplication with which is equivalent to the maps 
\[
\Omega_\ast \colon \GW(\C,\QF) \longrightarrow \GW(\C,\QF\qshift{-1-\sigma}) \quad \text{and} \quad \partial\colon \GW(\C,\QF) \longrightarrow \SS^1 \otimes \GW(\C,\QF\qshift{-1})
\]
for any Poincaré $\infty$-category $(\C;\QF)$; the former map here is induced by the equivalence $\Om\colon (\C,\QF) \rightarrow (\C,\QF\qshift{-1-\sigma})$ from Proposition~\refone{proposition:shiftofqingeneral}, and the latter is the boundary map of the Bott-Genauer sequence. The classes $\kappa$ and $\rho$ must then simply be given by applying these two maps to the unit of the multiplication, that is, to $[\mathbb S,q^\uni] \in \Omega^\infty\GW(\Spaf,\QF^\uni)$.
In the former case we thus set
\[
\kappa \coloneq \Omega_\ast[\mathbb S,q^\uni] = [\mathbb S^{-1},q^\uni] \in \Omega^\infty\GW(\Spaf,(\QF^\uni)\qshift{-1-\sigma}),
\]
and we will call $\kappa$ the \emph{Karoubi periodicity element}. In the latter case we put
\[
\delta \coloneq \partial[\mathbb S,q^\uni] \in \Omega^{\infty-1}\GW(\Spa^\fin,(\QF^\uni)\qshift{-1}),
\]
which we refer to as the \emph{Bott-Genauer element}. To give an explicit representative of it, recall that generally 
\[
\Omega^{\infty-1}\GW(\D,\Psi\qshift{-1}) \simeq |\Cob(\D,\Psi\qshift{-1})|,
\]
which receives a map from $\Poinc(\D,\Psi) \simeq \grpcr\Cob(\D,\Psi\qshift{-1})$. By Proposition \reftwo{proposition:boundarymapGB} this map agrees with the composite of $\partial$ and the tautological map $\Poinc(\C,\QF) \rightarrow \Omega^{\infty}\GW(\C,\QF)$, so that $\delta$ ends up simply represented by the object $(\mathbb S,q^\uni)$ of $\Cob(\Spaf,(\QF^\uni)\qshift{-1})$.



To see that multiplication with these classes does indeed agree with the indicated maps, note in the case of $\kappa$ that the Poincaré functor $\Omega \colon (\C,\QF) \rightarrow (\C,\QF\qshift{-1-\sigma})$ naturally identifies with 
\[
\id \otimes \Om\colon (\C,\QF) \otimes (\Spa^\fin,\QF^\uni) \longrightarrow (\C,\QF) \otimes (\Spa^\fin,(\QF^\uni)\qshift{-1-\sigma}).
\]
That the induced map $\Omega_\ast \colon \GW(\C,\QF) \longrightarrow \GW(\C,\QF\qshift{-1-\sigma})$ is given by multiplication with $\kappa$ is then witnessed by the lax symmetric monoidal structure of $\GW$, established in Corollary~\reffour{corollary:GW-universal-multiplicative}. More generally,
if $\F\colon \Catp \to \E$ is an arbitrary additive functor to some $\infty$-category $\E$ with finite limits then $\F$ is a $\GWspace$-module (for the induced $\Sps$-module structure on $\cP(\E) = \Fun(\E\op,\Sps)$ over $\Sps$, say) by Corollary~\reffour{corollary:GW-universal-multiplicative} and so for $(\C,\QF) \in \Catp$ we obtain a commutative square
\[
\begin{tikzcd}[column sep = 30pt]
\F(\C,\QF) \otimes \GWspace(\Spaf,\QF^{\uni}) \ar[d,"\mu"]\ar[r,"{(\id,\Om_*)}"] & \F(\C,\QF) \otimes \GWspace(\Spaf,(\QF^{\uni})\qshift{-1-\sig}) \ar[d,"\mu"] \\
\F(\C,\QF) \ar[r,"\Om_*"] & \F(\C,\QF\qshift{-1-\sig})
\end{tikzcd}
\]
in $\cP(\E)$, natural in $(\C,\QF)$, witnessing the fact that the lower horizontal map is induced by multiplication with $\kappa$. 
If, furthermore $\mathcal F$ is a module over some lax monoidal additive functor $\mathcal G \colon \Catp \rightarrow \Sps$, and $(\C,\QF)$ is a module over a monoidal Poincaré $\infty$-category $(\D,\Psi)$, then the action of $\GWspace(\Spaf,\QF\qshift{-1-\sigma})$ on $\F$ above factors through that of $\mathcal G(\D,\Psi\qshift{-1-\sigma})$, and so the action of $\kappa$ on $\F(\C,\QF)$ only depends on its image in $\mathcal G(\D,\Psi\qshift{-1-\sigma})$ under the unit map. We will continue to denote such images under unit maps by $\kappa$; for example, we obtain $\kappa \in \Lspace(\Dperf(\mathbb Z),\QF^{\mathrm g\sym})$ this way. 








The analogous statements for $\delta$ are obtained almost verbatim by noting that according to Lemma \refone{lemma:pairing-projection} the metabolic Poincaré-Verdier sequence
\[
(\C,\QF\qshift{-1}) \to \Met(\C,\QF) \xrightarrow{\met} (\C,\QF)
\]
of Example~\reftwo{example:metabolicfseq} is naturally obtained from that of $(\Spaf,\QF^{\uni})$ by tensoring with $(\C,\QF)$. Arguing as above, we then get that for any additive $\F\colon \Catp \to \E$ to a stable $\E$, the boundary map $\partial\colon \F(\C,\QF) \rightarrow \mathbb S^1 \otimes \F(\C,\QF\qshift{-1})$ is induced by multiplication with $\delta$. Again we will regard $\delta$ as an element of $\Omega^{\infty-1}\mathcal G(\D,\Psi)$ for any monoidal $(\D,\Psi) \in \Catp$ and lax monoidal additive $\mathcal G \colon \Catp \rightarrow \Spa$ without changing notation. 

Finally, note that as the Poincaré functors $\Omega$ are equivalences, the Karoubi periodicity element $\kappa$ is invertible in the sense that there exists a unique element $\kappa^{-1} \in \mathcal G(\D,\Psi\qshift{1+\sig})$ such that $\kappa\kappa^{-1}=\kappa^{-1}\kappa=1$. The same is of course generally true for the images of $\delta$ if and only if $\mathcal G$ is bordism invariant.  
In particular, the image $\delta \in \Om^{\infty-1}\L(\Spaf,(\QF^\uni)\qshift{-1})$ is invertible with inverse an element $\delta^{-1} \in \Om^{\infty+1}\L(\Spaf,(\QF^\uni)\qshift{1})$. %
We may therefore also consider
\[
\rho \coloneq \kappa \cdot \delta^{-2} \in \Om^{\infty+2}\L(\Spaf,(\QF^\uni)\qshift{1-\sigma})
\]
which we will refer to as the \emph{Ranicki periodicity element}. Multiplication by $\rho$ provides an equivalence
\[
\rho \cdot - \colon \F(\C,\QF) \longrightarrow \SS^{-2}\otimes \F(\C,\QF\qshift{1-\sigma})
\]
for every bordism invariant additive functor $\mathcal F$ which, by construction, agrees with the composite equivalence
\[
\F(\C,\QF) \xrightarrow{\partial^{-1} \circ \partial^{-1}} \SS^{-2}\otimes \F(\C,\QF\qshift{2}) \xrightarrow{\Omega_*} \SS^{-2}\otimes \F(\C,\QF\qshift{1-\sigma}).
\]
To compare this element to Ranicki's incarnation of it, let us finally give an explicit representative for $\rho$. To this end, recall that there is a natural map $\Poinc(\C,\QF\qshift{-n}) \rightarrow \Omega^{\infty+n}\L(\C,\QF)$, that is traditionally used to identify the homotopy groups of $\L$-spectra as bordism groups: Using $\Omega^\infty\L(\C,\QF) = |\Poinc \rho(\C,\QF)|$ this originates from the equivalence $(\C,\QF\qshift{-n}) \hookrightarrow \rho_n(\C,\QF)$ onto the subcategory with trivial boundary through the general map $\fib(X_n \rightarrow X^{\partial \Delta^n}) \rightarrow \Omega^n|X|$ %
where $X^- \colon \mathrm s\Sps\op \rightarrow \Sps$ is the limit preserving extension of a simplicial space $X$. The element $\delta$ is then the image of $(\SS^{-1},q^\uni) \in \Poinc(\Spaf,(\QF^\uni)\qshift{-1-\sigma})$ under this map (for $n=2$): Since this Poincaré form represents $\kappa$ the claim boils down to the composite
\[
\Poinc(\C,\QF\qshift{-n}) \rightarrow \Omega^{\infty+n}\L(\C,\QF) \xrightarrow{\partial^{n}}  \Omega^{\infty}\L(\C,\QF)
\]
being the tautological map. For $n=1$ this follows immediately from the analogous (finer) statement for Grothendieck-Witt spaces in Proposition~\reftwo{proposition:boundarymapGB} and the general case reduces to this by considering the embedding of $\rho_n(\C,\QF)$ into the $n$-fold metabolic Poincaré $\infty$-category in Lemma~\reftwo{lemma:rho_vs_power_set}. 
\end{construction}






\begin{comment}
Consider then its boundary map %
\[
\partial\colon \GW(\C,\QF) \longrightarrow \SS^1 \otimes \GW(\C,\QF\qshift{-1})
\]
Applying this again to $(\mathbb S,q^\uni)$, we obtain
\[
\delta \coloneq \partial(\mathbb S,q^\uni) \in \Omega^{\infty-1}\GW(\Spa^\fin,(\QF^\uni)\qshift{-1}),
\]
which we refer to as the \emph{Bott-Genauer element}. To give an explicit representative, recall that 
\[
\Omega^{\infty-1}\GW(\C,\QF\qshift{-1}) \simeq |\Cob(\C,\QF\qshift{-1})|,
\]
which receives a map from $\Poinc(\C,\QF) \simeq \grpcr\Cob(\C,\QF\qshift{-1})$. By Proposition \reftwo{proposition:boundarymapGB} this map agrees with the composite of $\partial$ and the inclusion $\Poinc(\C,\QF) \rightarrow \Omega^{\infty-1}\GW(\C,\QF)$, so that $\delta$ ends up simply represented by $(\mathbb S,q^\uni)$ itself.

Again, it follows from Lemma~\refone{lemma:pairing-projection} that the metabolic sequence for $(\C,\QF)$ is naturally obtained from that of $(\Spaf,\QF^{\uni})$ by tensoring with $(\C,\QF)$, and so as above for any additive $\F\colon \Catp \to \E$ to a stable $\E$, the boundary map $\partial\colon \F(\C,\QF) \rightarrow \mathbb S^1 \otimes \F(\C,\QF\qshift{-1})$ is induced by multiplication with $\delta$. Again we will regard $\delta$ as an element of $\mathcal G(\D,\Psi)$ for any monoidal $\mathcal G$ and $(\D,\Psi)$ without changing notation. 

Finally, while $\delta$ is not invertible, its image in $\Om^{\infty-1}\L(\Spaf,(\QF^\uni)\qshift{-1})$ (or any other bordism invariant additive functor) is, resulting in an element $\delta^{-1} \in \Om^{\infty+1}\L(\Spaf,(\QF^\uni)\qshift{1})$. %
We may therefore also consider
\[
\rho \coloneq \kappa \cdot \delta^{-2} \in \Om^{\infty+2}\L(\Spaf,(\QF^\uni)\qshift{1-\sigma})
\]
which we will refer to as the \emph{Ranicki periodicity element}. Multiplication by $\rho$ provides an equivalence
\[
\rho \cdot - \colon \F(\C,\QF) \longrightarrow \SS^{-2}\otimes \F(\C,\QF\qshift{1-\sigma})
\]
for every bordism invariant additive functor which, by construction, agrees with the composite equivalence
\[
\F(\C,\QF) \xrightarrow{\partial^{-1} \circ \partial^{-1}} \SS^{-2}\otimes \F(\C,\QF\qshift{2}) \xrightarrow{\Omega_*} \SS^{-2}\otimes \F(\C,\QF\qshift{1-\sigma}).
\]
To give a concrete description of $\rho$, recall that there are tautological maps $\Poinc(\D,\Psi\qshift{-n}) \rightarrow \Omega^{\infty+n}\L(\D,\Psi)$. We claim that $\rho$ is then simply the image of $(\mathbb S^{-1},q^\uni) \in \Poinc(\Spaf,(\QF^\uni)\qshift{-1-\sigma})$ with $n=2$; this is tautological in the description of $\L(\D,\Phi)$ as the stabilisation of $\GW(\C,\QF)$, i.e. the colimit along the boundary map of the metabolic fibre sequence, as it simply identifies the tautological map above as the composite
\[
\Poinc(\D,\Psi\qshift{-2}) \longrightarrow \Omega^\infty\GW(\D,\Psi\qshift{-2}) \xrightarrow{\bord} \Omega^\infty\L(\D,\Psi\qshift{-2})\xrightarrow{\partial^2} \Omega^{\infty+2}\L(\D,\Psi).
\]
\end{construction}
\end{comment}










\begin{remark}\
\begin{enumerate}
\item For additive $\F$, we deduce from Proposition~\reftwo{proposition:Q-bordism} the formula 
\[
\F^\bord(\C,\QF) \simeq \colim [\F(\C,\QF) \xrightarrow{\delta\cdot -} \Sigma \F(\C,\QF\qshift{-1}) \xrightarrow{\delta\cdot -} \Sigma^2 \F(\C,\QF\qshift{-2}) \to \cdots ]
\]
and thus $\F^\bord = \F[1/\delta]$, where $\F[1/\delta]$ refers to the initial functor under $\F$ for which multiplication by $\delta$ is an equivalence, or equivalently, the initial functor under $\F$ which vanishes on $\Met(\C,\QF)$ for all Poincar\'e $\infty$-categories; that this initial functor agrees with $\F^\bord$ is a consequence of Lemma~\reftwo{lemma:criterion}. In particular, we have $\L = \GW[1/\delta]$; the latter term (defined as the above displayed filtered colimit) was used as a definition of L-theory in Schlichting's work, see \cite{schlichting-derived}*{Definition 7.1} (where the class we call $\delta$ corresponds to the class called $\eta$ in loc.\ cit.).
\item Recall from \S\reftwo{subsection:Mackey} that any additive $\F\colon \Catp \to \Spa$ has a canonically associated genuine refinement $\F^{\g\hyp} \colon \Catp \to \Spa^{\g\Ct}$ along the fixed point functor $(-)^{\g\Ct}$ such that the geometric fixed points of $\F^{\g\hyp}$ canonically identify with $\F^\bord$. The above formula $\F^\bord(\C,\QF) = \F(\C,\QF)[1/\delta]$ is then well-known in equivariant homotopy theory, where the geometric fixed points are obtained from the genuine fixed points by inverting the equivariant Euler class of the reduced regular representation: And indeed, unravelling the definitions, (and using that the reduced regular representation for $\Ct$ is precisely the sign representation) this equivariant Euler class for $\KR(\C,\QF)=\GW^{\g\hyp}(\C,\QF)$ is precisely the Bott-Genauer class $\delta$.
\end{enumerate}
\end{remark}






Let us now explain the relation of the Karoubi periodicity element $\kappa$ to his fundamental theorem. Recall that its general implementation in \papertwo is the composite
\[
\mathbb S^1 \otimes \mathrm V(\C,\QF) \simeq \GW(\C,\QF\qshift{1}) \simeq \mathrm U(\C,\QF\qshift{2})
\]
with both intermediate equivalences arising from the Bott-Genauer sequence, see Corollary~\reftwo{corollary:karoubi-fundamental-theorem}. Further composing with multiplication by $\kappa$ on the additive functor $\mathrm{U}$, we obtain an equivalence 
\[
\mathbb S^1 \otimes \mathrm V(\C,\QF) \simeq \mathrm U(\C,\QF\qshift{2}) \xrightarrow{\kappa \cdot -} \mathrm U(\C,\QF\qshift{1-\sigma}).
\]
which by construction agrees with the map used in \papertwo to establish our extensions of Karoubi's fundamental theorem, see the discussion after Corollary~\reftwo{corollary:karoubi-fundamental-theorem}. 



In the special case of discrete rings we can unwind further:

\begin{example}
\label{example:perclasses}%
Let $R$ be a commutative ring, $A$ and $R$-algebra and $M$ an invertible $R$-module with involution over $A$. Recall that in this case for a Poincaré structure on $\QF_M$ on $\Dperf(A)$, whose bilinear part is classified by $M$, the bilinear part of $\QF\qshift{1-\sigma}$ is classified by $-M$, simply given by inserting a sign into the original action. 
\begin{enumerate}
\item In the simplest case of the symmetric and quadratic structures we consequently have $(\QF^\sym_M)\qshift{1-\sigma} = \QF^\sym_{-M}$ and periodicity elements 
\[
\kappa \in \Om^{\infty}\GW(\Dfree(R),(\QF^{-\sym})\qshift{-2}) \quad \text{and} \quad \rho \in \Om^{\infty+2}\L(\Dfree(R),\QF^{-\sym}),
\]
both represented by the complex $R[-1]$ with its tautological 2-shifted skew-symmetric form. Multiplication with these elements now gives equivalences
\[
\mathbb S^1 \otimes \mathrm V(\D^c(A),\QF^r_{M}) \xrightarrow{\kappa \cdot -} \mathrm U(\D^c(A),\QF^r_{-M}) \quad \text{and} \quad \L(\D^c(A),\QF^r_M) \xrightarrow{\rho \cdot -} \SS^{-2}\otimes \L(\D^c(A),\QF^r_{-M}),
\]
respectively, where $r \in \{\qdr,\sym\}$ and $c \subseteq \K_0(A)$ is some choice of control group. By the discussion above these agree with the ones established in Corollaries~\reftwo{corollary:karoubi-fundamental} and~\reftwo{corollary:Ranickipergeneral}, respectively, 
and multiplication with $\rho^2 \in \Om^{\infty+4}\L(\Dfree(R),\QF^\sym)$, represented by $R[-2]$ with its tautological 4-shifted symmetric form 
gives Ranicki's classical periodicity isomorphisms
\[
\L(\D^c(A),\QF^r_M) \longrightarrow \SS^{-4}\otimes \L(\D^c(A),\QF^r_{M}).
\]
\item As for the genuine symmetric structure, we have $(\QF^{\mathrm g\sym}_R)\qshift{1-\sigma} \simeq \QF^{-\gev}_R$ by Corollary~\refone{corollary:shift-genuine}, and so we may view the Karoubi and Ranicki periodicity elements as living in
\[
\kappa \in \Om^{\infty}\GW(\Dfree(R),(\QF^{-\gev})\qshift{-2}) \quad \text{and} \quad \rho \in \Om^{\infty+2}\L(\Dfree(R),\QF^{-\gev}).
\]
If $\lambda$ is an invertible form parameter of symmetric type on $A$ with underlying $R$-module with involution $M$, then $(\QF^{\mathrm g\lambda}_M)\qshift{1-\sigma} = \QF^{\mathrm g \check{\lambda}}_{-M}$, see Proposition~\refone{proposition:form-parameter-shift}, and we find that multiplication with the periodicity classes
gives the equivalences
\[
\mathbb S^1 \otimes \mathrm V(\D^c(A),\QF^{\mathrm g\lambda}_{M}) \longrightarrow \mathrm U(\D^c(A),\QF^{\mathrm g\check{\lambda}}_{-M}) \quad \text{and} \quad \mathbb S^{2} \otimes \L(\D^c(A),\QF^{\mathrm g\lambda}_{M}) \longrightarrow \L(\D^c(A),\QF^{\mathrm g\check{\lambda}}_{-M}),
\]
from Corollary~\reftwo{corollary:formparakarper} and%
~\reftwo{corollary:karoubi-periodicity-form-parameters}, the former of which is very much in the spirit of Karoubi's original formulation in \cite{karoubi-periodicity}. 

\end{enumerate}
\end{example}

\begin{remark}
To avoid confusion, let us mention the following misalignment of the example above with \cite{karoubi-periodicity}. Karoubi predicted the existence of an element $x \in \Omega^{\infty+2}\GW^{-\mathrm g\sym}(\mathbb Z)$ analogous to Ranicki's periodicity element $\rho \in \Omega^{\infty+2}\L^{-\mathrm g\sym}(\mathbb Z)$, multiplication with which ought to induce an equivalence
\[
\mathbb S^1 \otimes \mathrm V(\D^c(A),\QF^{\mathrm g\lambda}_{M})\longrightarrow\mathrm U(\D^c(A),\QF^{\mathrm g\check{\lambda}}_{-M}),
\]
in place of our $\kappa \in \Omega^\infty\GW(\Dfree(\mathbb Z),(\QF^{-\mathrm g\ev})\qshift{-2}) = \Omega^\infty\GW(\Dperf(\mathbb Z),(\QF^{-\mathrm g\sym})\qshift{-2})$. The only multiplication of the above shape induced by such a class $x$ seems to be that with $\partial^2(x) \in \Omega^\infty\GW(\Dperf(\mathbb Z),(\QF^{-\gs})\qshift{-2})$. However, any element of $\GW_0(\Dperf(\mathbb Z),(\QF^{-\gs})\qshift{-2})$ that lies in the image of $\partial$ vanishes in $\K_0(\mathbb Z)$. Thus multiplication with it kills all hyperbolic forms and it consequently cannot induce an equivalence as above for all $A$. 

In fact, in Proposition~\refthree{proposition:gwz-low-deg} we calculated $\GW_2^{-\mathrm g\sym}(\mathbb Z) \cong \mathbb Z$ and from the fundamental fibre sequence one reads off that the bordification map to $\L_2^{-\mathrm g\sym}(\mathbb Z) = \mathbb Z$ is injective with index $4$. Similarly one finds 
\[
\GW_0(\Dperf(\mathbb Z),(\QF^{-\mathrm g\sym})\qshift{-2}) \cong \mathbb Z \times_{\mathbb Z/2} \mathbb Z
\]
via signature and Euler characteristic. It follows that $\partial^2$ takes classes in $\GW_2^{-\mathrm g\sym}(\mathbb Z)$ to multiples of $(4,0)$, while $\kappa$ corresponds to $(1,-1)$, and also that Ranicki's periodicity class $[\rho] \in \L_2^{-\g\sym}(\mathbb Z)$ does not lift to $\GW_2^{-\g\sym}(\mathbb Z)$. 

In total we consider the above implementation of the equivalence in Karoubi's conjecture as the closest possible approximation to the original formulation.
\end{remark}




Let us now extend the discussion of Ranicki's periodicity isomorphism 
\[
\rho^2 \cdot -\colon  \L(\D^c(R),\QF^\sym_M) \longrightarrow \SS^{-4}\otimes \L(\D^c(R),\QF^\sym_{M}).
\]
and its relatives to $n\sigma$-oriented $\Einf$-rings $R$ for $n \geq 1$, in place of ordinary rings (where $n=2$). 
To this end, recall from Definition~\refone{def:nsigmaorpaperI} that an $n\sigma$-orientation on an $\Einf$-ring with involution $(R,\tau)$ is a
$\Ct$-equivariant $R$-linear equivalence $\eta\colon R \rightarrow \mathbb S^{n\sigma-n} \otimes R$ %
lifting the identity on underlying $R$-modules, where $\Ct$ acts on the left term via $\tau$ on the right term via the corresponding tensor of actions. %
Equivalently, the data of such an equivalence is a lift 
\[
\eta\in (\QF^{\sym}_{\tau})\qshift{n\sigma-n}(R) = \Omega^\infty (\mathbb S^{n\sigma-n} \otimes R)^\hC
\]
of the unit $1 \in \Omega^\infty R$ along the forgetful map $\Omega^\infty (\mathbb S^{n\sigma-n} \otimes R)^\hC \rightarrow \Omega^\infty R$. Note that if $R$ is discrete, then it indeed carries a unique $2\sigma$-orientation, since this forgetful map is then an equivalence. 

In fact, it is convenient to upgrade the notion of orientation to general Poincaré $\infty$-categories:

\begin{definition}
For a natural number $n \geq 1$, an $n\sigma$-orientation on a commutative algebra $(\D,\Phi)$ in $\Catp$ is an element $\eta \in \Omega^\infty\Phi\qshift{n\sigma-n}(\mathbb 1)$ lifting the unit of
\[
\Omega^\infty \Bil_{\Phi\qshift{n\sigma-n}}(\mathbb 1, \mathbb 1) \simeq \Omega^\infty \Bil_{\Phi}(\mathbb 1, \mathbb 1)
\]
under the forgetful map $\Phi\qshift{n\sigma-n} \Rightarrow \Bil_{\Phi\qshift{n\sigma-n}} \circ \Delta$. 
\end{definition}

An $n\sigma$-orientation on an $\Einf$-ring $R$ with involution $\tau$ is now by definition the same thing as an $n\sigma$-orientation on $(\Dperf(R),\QF^\sym_\tau)$.   


\begin{construction}
An $n\sigma$-orientation $\eta$ on $(\D,\Phi)$ determines a Poincaré object $o_\eta = (\one,\eta) \in \Poinc(\D,\Phi\qshift{n\sigma-n})$, multiplication with which gives Poincaré functors
\[
(\D,\mathbb \Phi\Phi\qshift{-n\sigma-n}) \longrightarrow (\D,\Phi\qshift{-2n}) \quad \text{and} \quad (\D,\mathbb \Phi\Phi\qshift{-n\sigma+n}) \longrightarrow (\D,\Phi)
\]
covering the identity of $\D$, and thus
\[
o_{\eta}\kappa^n \in \Om^{\infty}\GW(\D,\Phi\qshift{-2n}) \quad \text{and} \quad o_{\eta}\rho^n \in \Om^{\infty+2n}\L(\D,\Phi),
\]
which we will refer to as the \emph{potential periodicity elements} of $\eta$. For any module $(\C,\QF)$ over $(\D,\Phi)$ they induce maps
\[
o_{\eta}\kappa^n \cdot - \colon \GW(\C,\QF) \longrightarrow \GW(\C,\QF\qshift{-2n}) \quad \text{and} \quad  o_{\eta}\rho^n \cdot - \colon \L(\C,\QF) \longrightarrow \SS^{-2n}\otimes \L(\C,\QF).
\]
\end{construction}



\begin{comment}
\begin{construction}
Viewing $\eta$ as an equivalence of $(A,\tau)$-modules in $\Mod_{A\otimes A}^{\hC}=\Funs(\Modp{A})$, 
it determines 
a refinement of the identity on $\Modp{A}$ to a $(\Modp{A},\QF^{\sym}_{\tau})$-linear equivalence 
\[
(\Modp{A},\QF^{\sym}_{\tau}) \simeq (\Modp{A},(\QF^{\sym}_{\tau})\qshift{n\sigma-n}) ,
\]
which we may consequently describe as multiplication with the Poincaré object 
\[
o_{\eta} = (A,q_{\eta}) \in \Poinc(\Modp{A},(\QF^{\sym}_{\tau})\qshift{n\sigma-n}).
\]
By abuse of notation let us also write $o_{\eta}$ for the associated point in $\Om^{\infty}\GW(\Dperf(R),(\QF^{\sym}_{\tau})\qshift{n\sigma-n})$.
Since $\L$ is lax symmetric monoidal by Corollary~\reffour{corollary:lmonoidal},
we then conclude that Ranicki periodicity for $A$, given by definition as the composite of equivalences
\[
\begin{tikzcd}
\L(\Modp{A},\QF^{r}_{\tau}) \ar[r,"\rho^n\cdot-","\simeq"'] & \SS^{-2n}\otimes \L(\Modp{A},(\QF^{r}_{\tau})\qshift{n-n\sig}) \ar[r,"o_{\eta}\cdot-","\simeq"'] & \SS^{-2}\otimes \L(\Modp{A},\QF^{r}_{\tau})
\end{tikzcd}
\]
for $r \in \{\qdr,\sym\}$, is therefore itself multiplication with $o_{\eta}\rho^n\in \Om^{\infty-2n}\L(\Modp{A},\QF^{\sym}_{\tau})$. This element is then graded-invertible (with inverse lying in 
$\Om^{\infty+2}\L(\Modp{A},\QF^{\sym}_{\tau})$). Similarly, if $A$ is connective then by Corollary~\refone{corollary:periodicity-poincare} an $n\sigma$-orientation $\eta$ determines an $(\Modp{A},\QF^{\geq 0}_{\tau})$-linear refinement
\[
(\Modp{A},\QF^{\geq 0}_{\tau}) \simeq (\Modp{A},(\QF^{\geq n}_{\tau})\qshift{n\sig-n})
\]
of the identity, which we may identify as multiplication with the Poincaré object 
\[
o^{\g}_{\eta} = (A,q_{\eta}) \in \Poinc(\Modp{A},(\QF^{\geq n}_{\tau})\qshift{n\sigma-n}),
\]
where $q_{\eta}$ is interpreted as a form via the equivalence $\Om^{\infty}((\QF^{\geq n}_{\tau})\qshift{n\sigma-n}(A)) \xrightarrow{\simeq} \Om^{\infty}((\QF^{\sym}_{\tau})\qshift{n\sigma-n}(A))$ induced by polarization (in light of Corollary~\refone{corollary:periodicity-poincare}).
Multiplication by $o^{\g}_{\eta}\rho^n\in \Om^{\infty-2n}\L(\Modp{A},\QF^{\geq n}_{\tau})$ then gives the connective Ranicki periodicity equivalences
\[
\L(\Modp{A},\QF^{\geq m}_{\tau}) \to \SS^{-2n}\otimes \L(\Modp{A},\QF^{\geq m+n}_{\tau})
\]
for $m \in \ZZ$. 
\end{construction}

We note that the case of discrete rings described above relies on the unique $2\sig$-orientation enjoyed by all discrete rings, though given its uniqueness we have allowed ourselves to omit it from the notation when discussed there.

To account also for non-invertible periodicity maps it will be convenient to introduce the following extension of the notion of orientation to general Poincaré $\infty$-categories.
\end{comment}


 



\begin{examples}\
\begin{enumerate}
\item
For an $\Einf$-ring spectrum with involution $(R,\tau)$ an $n\sig$-orientation $\eta$ on $(\Modf{R},\QF^{\sym}_{\tau})$ amounts to an $n\sigma$-orientation on $(R,\tau)$. In this case the Poincaré object $o_{\eta}$ is invertible (since for symmetric Poincaré $\infty$-categories this is detected on underlying objects) and the composite equivalence
\[
(\Modf{R},\QF^{\sym}_{\tau}) \xrightarrow{\Om^n} (\Modf{R},(\QF^{\sym}_{\tau})\qshift{-n-n\sig}) \xrightarrow{o_{\eta}} (\Modf{R},(\QF^{\sym}_{\tau})\qshift{-2n})
\]
coincides (up to a shift) with the equivalence of Corollary~\refone{corollary:periodicity-poincare} specialised to the case at hand, since $\eta\colon R \rightarrow \mathbb S^{n\sigma-n} \otimes R$ is $R$-linear.
In particular, the potential periodicity elements 
\[
o_{\eta}\kappa^n \in \Om^{\infty}\GW(\Modf{R},(\QF^\sym_\tau)\qshift{-2n}) \quad \text{and} \quad o_{\eta}\rho^n \in \Om^{\infty+2n}\L(\Modf{R},\QF_\tau^\sym)
\]
are invertible, 
and we find periodicity equivalences 
\[
o_\eta\rho^n \cdot -\colon  \L(\Mod^c(R),\QF^r_M) \longrightarrow \SS^{-2n}\otimes \L(\Mod^c(R),\QF^r_{M})
\]
for $r \in \{\mathrm q,\sym\}$ recovering the ones from Corollary~\reftwo{corollary:Ranickipergeneral}. The same is true for any $\Eone$-$R$-algebra $A$ with invertible $R$-module with involution.

A completely analogous statement also holds, whenever $\D$ is any symmetric monoidal $\infty$-category and $\Phi = \QF^{\sym}_{\Dual}$ is the symmetric Poincaré structure associated to some symmetric monoidal duality on $\D$ (that is, a duality yielding a commutative algebra object in $\Catps$), and $(\C,\QF)$ is a module over $(\D,\Psi)$ for which $\QF$ is either symmetric or quadratic. 
\item
If $(R,\tau)$ is a connective $\Einf$-ring spectrum with involution then an $n\sigma$-orientation uniquely lifts along the map $\QF_\tau^{\geq 0} \rightarrow \QF_\tau^\sym$, since the map
\[
\Omega^\infty (\QF_{\tau}^{\geq 0})\qshift{n\sigma-n}(R) \longrightarrow \Omega^\infty (\QF_{\tau}^{\sym})\qshift{n\sigma-n}(R)
\]
is then an equivalence. The potential periodicity elements 
\[
o_{\eta}\kappa^n \in \Om^{\infty}\GW(\Modf{R},(\QF^{\geq 0}_\tau)\qshift{-2n}) \quad \text{and} \quad o_{\eta}\rho^n \in \Om^{\infty+2n}\L(\Modf{R},\QF_\tau^{\geq 0})
\]
are then generally not invertible, the latter giving rise to the stabilisation maps
\[
o_{\eta}\rho^n \cdot - \colon \L(\Mod^c(R),\QF^{\geq k}_{\tau}) \to \SS^{-2n}\otimes \L(\Mod^c(R),\QF^{\geq k}_{\tau}).
\]
for arbitrary $k \in \mathbb Z \cup \{\pm \infty\}$. 

\item In the situation of the previous point, the polarisation map
\[
\Omega^\infty(\QF^{\geq n}_\tau)\qshift{n\sigma-n}(R) \longrightarrow \Omega^\infty(\QF^{\geq 0}_\tau)\qshift{n\sigma-n}(R)
\]
is an equivalence, yielding a lift $o^\mathrm g_\eta \in \Poinc(\Modf{R},(\QF^{\geq n}_\tau)\qshift{n\sigma-n})$ of $o_\eta$ and consequently a factorisation
\[
(\Modf{R},(\QF^{\geq 0}_\tau)\qshift{-n\sigma+n}) \xrightarrow{o_\eta^g \cdot -} (\Modf{R},\QF^{\geq n}_\tau) \xrightarrow{\mathrm{pol}} (\Modf{R},\QF^{\geq 0}_\tau),
\]
of the multiplication with $o_\eta$, whose first map is an equivalence by direct inspection of linear parts. Multiplication with $o_\eta^\g$ then recovers the equivalences 
\[
o^\mathrm g_{\eta}\rho^n \cdot - \colon \L(\Mod^c(R),\QF^{\geq k}_{\tau}) \to \SS^{-2n}\otimes \L(\Mod^c(R),\QF^{\geq k+n}_{\tau}).
\]
induced by the equivalences of Corollary~\refone{corollary:periodicity-poincare-truncated}, refining the stabilisation maps from the previous point.
\item 
The simplest example that is not quite of this nature 
is $\Dperf(\ZZ)$ equipped with the Poincar\'e structure $\QF^{\g\fr}_{\ZZ}$ associated to the Frobenius
form parameter of Example~\reffour{example:form_parameter_over_Tate_Poincare}.
In this case $(\QF^{\g\fr}_{\ZZ})\qshift{2\sigma-2}(\ZZ)$ evaluates to the (Eilenberg-Mac Lane spectrum of the) pullback of
\[
\ZZ^\hC \longrightarrow \ZZ^\tC \longleftarrow \ZZ[-2] \times_{\ZZ/2[-2]} \tau_{\geq -2}\ZZ^\tC.
\]
In particular, at the level of infinite loop spaces, the fibre over the unit of $\Bil_{\QF^{\g\fr}_{\ZZ}}(\ZZ, \ZZ)^{\hC} \simeq \ZZ^\hC$ is trivial, giving a unique $2\sigma$-orientation $\eta$ to $\QF^{\g\fr}_{\ZZ}$ lifting the $2\sigma$-orientation of $\QF^{\sym}_{\ZZ}$. The associated potential periodicity element $o_{\eta}\rho^2  \in \Om^{\infty}\L^{\g\fr}(\ZZ)$ then lifts the one in $\Om^{\infty+4}\L^{\sym}(\ZZ)$; in fact the map between these two $\Om^{\infty+4}\L$ spaces is an equivalence by Example~\refthree{example:dotto-ogle}, and so such a potential periodicity element is unique as well.
Since for every ring with form parameter $(R,M,\lambda)$, the spectrum $\L(\D^c(R),\QF^{\mathrm g\lambda}_M)$ is a module over $\L^{\g\fr}(\ZZ)$
by Proposition~\reffour{prop:tensormackeyvstensorfpm}, we learn that these spectra always carry stabilisation maps
\[
\L(\D^c(R),\QF^{\mathrm g\lambda}_M) \longrightarrow \SS^{-4}\otimes \L(\D^c(R),\QF^{\mathrm g\lambda}_M)
\]
induced by the multiplication with the potential periodicity element $o_{\eta}\rho^2 \in \Om^{\infty+4}\L^{\g\fr}(\ZZ)$ 
(and by its image in $\Om^{\infty+4}\L^{\mathrm g\sym}(\mathbb Z)$ when $\lambda$ is of symmetric type). 
\item 
Not every lax-symmetric monoidal Poincar\'e structure over an $n\sigma$-oriented $\Einf$-ring admits an $n\sigma$-orientation itself: The simplest example is the Tate Poincar\'e structure on $\ZZ$: In this case
$(\QF^\tate_{\ZZ})\qshift{2\sigma-2}(\ZZ)$ evaluates to the pullback of
\[
\ZZ^\hC \longrightarrow \ZZ^\tC \longleftarrow \ZZ[-2],
\]
which does not admit the required lift of the element $1 \in \Om^{\infty}(\ZZ^\hC)$. And indeed, the forthcoming calculation of $\L^{\tate}(\ZZ)$ in \cite{TateL} show that Ranicki's periodicity element in $\Om^{\infty+4}\L^{\g\fr}(\ZZ)$ does not admit a further lift to this spectrum.
\end{enumerate}
\end{examples}


























\subsection{Multiplicativity of the circle decomposition theorem}
\label{subsection:mult-shaneson}%








Let $\F\colon \Catp \to \A$ be a functor valued in a cocomplete $\infty$-category $\A$ %
and recall from \S\reffour{subsection:shaneson} the assembly map
\[
\B\ZZ \otimes \F(\C,\QF) \to \F((\C,\QF)_{\B\ZZ})
\]
defined for any Poincaré $\infty$-category $(\C,\QF)$ as the colimit-exchange map. In this section we refine this map to a symmetric monoidal transformation in the case where $\A$ is symmetric monoidal and $\F$ is lax symmetric monoidal. We begin with the following general observation.


\begin{proposition}
Let $\A$ be a cocomplete symmetric monoidal $\infty$-category in which the tensor product preserves colimits in each variable, $\F\colon \Catp \to \A$ a lax symmetric monoidal functor, $\I$ a small symmetric monoidal $\infty$-category and $\varphi \colon \I \to \Catp$ a functor. Then the associated assembly map
\[
\colim_\I \F\circ \varphi \to \F(\colim_\I \varphi)
\]
is the component of a lax symmetric monoidal transformation among lax symmetric monoidal $\A$-valued functors in $\varphi$. In particular, if $\varphi$ is lax symmetric monoidal, then the map 
\[
\colim_\I \F\circ \varphi \to \F(\colim_\I \varphi)
\]
is one of commutative algebra objects in $\A$.
\end{proposition}
\begin{proof}
Recall from Remark~\reffour{remark:left-kan-monoidal} that under the assumptions of the proposition the colimit functor refines to a symmetric monoidal adjunction
\[
\colim_\I \colon \Fun(\I,\Catp)^{\Day} \adj \Catp \cocolon \const
\]
where the symmetric monoidal structure on the left is given by Day convolution, %
$\colim_\I$ is symmetric monoidal, %
and $\const$, the constant diagram functor, is canonically lax symmetric monoidal. In particular, the associated unit and counit maps are lax symmetric monoidal transformations among lax symmetric monoidal functors.
Since $\F$ is lax symmetric monoidal,
the associated commutative square 
\[
\begin{tikzcd}
\Catp \ar[r,"\F"]\ar[d,"\const"'] & \A \ar[d,"\const"] \\
\Fun(\I,\Catp)^{\Day} \ar[r,"\F_*"] & \Fun(\I,\A)^{\Day}
\end{tikzcd}
\]
of lax symmetric monoidal functors determines a (Beck-Chevalley type) natural transformation
\[
\colim_\I \circ \F_* \Rightarrow \colim_\I \circ\F_*\circ \const\circ\colim_\I = \colim_\I\circ \const\circ \F \circ\colim_\I \Rightarrow \F\circ\colim_\I
\]
which, being a composite of lax symmetric monoidal transformations, gives the desired lax symmetric monoidal transformation.
\end{proof}


Combining this with Corollary~\reffour{corollary:iterated-shaneson-splitting} we conclude:
\begin{corollary}
\label{corollary:iteratedassmonoidal}%
Let $\A$ be a cocomplete symmetric monoidal $\infty$-category in which the tensor product preserves colimits in each variable and let $\F\colon \Catp \to \A$ be a lax symmetric monoidal bordism invariant Karoubi-localising functor. %
Then for every $n \geq 0$ the assembly map
\[
\B\ZZ^n \otimes \F(-) \Rightarrow \F((-)_{\B\ZZ^n})
\]
is an equivalence of lax symmetric monoidal functors $\Catp \to \A$. %
\end{corollary}


\begin{example}
\label{example:group-rings}%
For $(R,\sig)$ a commutative ring with involution the symmetric monoidal Poincaré $\infty$-category $(\Dperf(R),\QF^{\sym}_{\sig})_{\B\ZZ^n}$ is equivalent to $(\Dperf(R[\ZZ^n]),\QF^{\sym}_{\bar{\sig}})$ where we view $R[\ZZ^n]$ as a ring with involution $\bar{\sig}$ extending $\sig$ and inducing the inversion on $\ZZ^n$. By Corollary~\reffour{corollary:iteratedassmonoidal} we then obtain 
a canonical equivalence of $\Einf$-rings
\[
\KL^{\sym}(R[\ZZ^n],\bar{\sig}) \simeq  \KL^{\sym}(R,\sig)[\B\ZZ^n].
\]
where for an $\Einf$-ring $A$ we denote by
$A[\B\ZZ^n]=\B\ZZ^n \otimes A$ the associated group-algebra.
\end{example}

For not necessarily bordism invariant Karoubi-localising lax symmetric monoidal functors $\F$, like $\KK$ or $\KGW$, the map 
\[
\B\ZZ^n \otimes \F(-) \Rightarrow \F((-)_{\B\ZZ^n})
\]
is still one of lax symmetric monoidal functors, and an equivalence in certain cases:

\begin{example}
\label{example:regular-noetherian}%
For $(R,\sig)$ a regular noetherian commutative ring with involution 
the map of $\Einf$-rings
\[
\KGW^{\sym}(R,\sig)[\B\ZZ^n] \to \KGW^{\sym}(R[\ZZ^n],\bar{\sig})
\]
induced by the assembly map, is an equivalence. Indeed, in these cases the inductively appearing complementary summands $\N\KK(\Dperf(R[\ZZ^i]))$ vanish. 
\end{example}











Let us now specialise to the case $n=1$, for which Examples \reffour{example:group-rings} and \reffour{example:regular-noetherian} motivate one to understand the $\Einf$-algebra structure on 
the group $\Einf$-$A$-algebra $A[\B\ZZ]$ associated to an $\Einf$-ring spectrum $A$ and the $\Einf$-group $\B\ZZ$. Being a group algebra, it comes equipped with an augmentation $A[\B\ZZ] \to A$ induced by the terminal map $\B\ZZ \to \ast$ whose augmentation ideal is given by $\Sigma A$. Hence $A[\B\ZZ]$ is equivalent to $A\oplus\Sigma A$ as an $A$-module over and under $A$, and one may ask if this equivalence is multiplicative when viewing $A \oplus \Sigma A$ as
the split square-zero extension.


\begin{lemma}
\label{lemma:sq0}%
Let $A$ be an $\Einf$-algebra. Refinements of the $A$-module-over-under-$A$ equivalence $A[\B\ZZ] \simeq A \oplus \Sigma A$ to one of augmented  
$\Ek{n}$-$A$-algebras naturally correspond to factorisations of the unit $\SS \to A$,
as a map of spectra, 
through the map $\SS \to \SS^{-n-1}\otimes\Sig^{\infty}K(\ZZ,n+1)$. In particular, for $n=\infty$, such refinements correspond to factorisations of the unit though the Eilenberg-MacLane spectrum $\GEM \ZZ = \colim_n\SS^{-n-1}\otimes\Sig^{\infty}K(\ZZ,n+1)$.
\end{lemma}
\begin{proof}
For an $A$-module $M$, we claim that $\fib(\gl_1(A\oplus M) \to \gl_1(A))$ is naturally equivalent to $\Omega^\infty M$: Indeed, as functors of $M$ both expressions are right adjoints $\Mod_A \to \Grp_{\Einf}$ among pointed presentable $\infty$-categories which agree after post-composition with the right adjoint $\Grp_{\Einf} \to \Sps_*$, and hence must agree before post-composition since $\Mod_A$ is stable and the induced functor $\Sp(\Grp_{\Einf}) \to \Sp(\Sps_*)$ is an equivalence.
Given this claim, the desired result now follows from the following calculation of the mapping space of augmented $\Ek{n}$-$A$-algebras:
\begin{align*}
\Map_{\Alg^{\mathrm{aug}}_{\Ek{n}}(\Mod_A)}(A[\B\ZZ],A\oplus \Sig A) &= \Map_{\Grp_{\Ek{n}}}(\B\ZZ,\fib[\gl_1(A\oplus \Sig A) \to \gl_1(A)]) \\
&= \Map_{\Grp_{\Ek{n}}}(\B\ZZ,\Om^{\infty}(\Sigma A))\\
&= \Map_{\Grp_{\Ek{n+1}}}(\ZZ,\Om^\infty A) \\
&= \Map_{\Sp}(\Om^{n+1}\Sig^{\infty}K(\ZZ,n+1),A).\qedhere
\end{align*}
\end{proof}

As a sample application of the above, we obtain:
\begin{corollary}
Let $(R,\sig)$ be a discrete commutative ring with involution. Then there is a preferred equivalence
\[
\KL^\s(R[\ZZ],\bar{\sig}) \simeq \KL^\s(R,\sig) \oplus \Sig \KL^\s(R,\sig)
\]
as augmented $\Eone$-$\KL^\s(R,\sig)$-algebras which in general does not extend to one of $\Ek{2}$-algebras. After localisation at 2, however, there is such an equivalence as $\Einf$-algebras.
\end{corollary}
\begin{proof}
Note that for $n=1$, the condition of Lemma~\reffour{lemma:sq0} is equivalent to equipping $A$ with a complex orientation and recall that $\L^\s(\ZZ)$ admits a preferred complex orientation, the Sullivan-Ranicki orientation. In particular, from the $\Einf$-map $\L^\s(\ZZ) \to \KL^\s(R,\sig)$, we deduce that $\KL^\s(R,\sig)$ is equipped with an induced complex orientation, so that the first claim follows from combining Example~\reffour{example:group-rings} with Lemma~\reffour{lemma:sq0} in case $n=1$. Then note that for $n=\infty$, the condition of Lemma~\reffour{lemma:sq0} holds whenever the underlying spectrum of $A$ admits a $\ZZ$-module structure. This is the case for $\KL^\s(R,\sig)_{(2)}$ as is again transported from the Sullivan-Ranicki orientation and the classical fact that $\MSO_{(2)}$ admits a $\ZZ$-module structure. Finally, to see the claim about $\Ek{2}$-structures, we show that $\L^\s(\ZZ)[\B\ZZ]$ is not $\Ek{2}$-equivalent to the trivial square-zero extension $\L^\s(\ZZ) \oplus \Sig \L^\s(\ZZ)$. By Lemma~\reffour{lemma:sq0}, it suffices to show that there is no map $K(\ZZ,3) \to \Omega^\infty \Sigma^3\L^\s(\ZZ)$ inducing an isomorphism on $\pi_3$. If such a map exists, we conclude from the equivalence $\L^\s(\ZZ)[1/2] \simeq \KO[1/2]$ that $K(\ZZ,3)$ is, after localising away from 2, a retract of $\Omega^\infty \Sigma^3 \KO[1/2]$ which in turn is a retract of $\Omega^\infty \Sigma^3 \ku[1/2] \simeq \SU[1/2]$. But the $\ZZ$-homology of $\SU[1/2]$ is torsion-free, while $H_8(K(\ZZ,3),\ZZ)$ contains non-trivial $3$-torsion.
\end{proof}

\begin{example}
\label{remark:L-group-sq0}%
As a consequence of the above, for a commutative ring $(R,\sig)$ with involution, $\KL^{\sym}_0(R,\sig)[\B\ZZ]$ is a square-zero extension of $\KL^{\sym}_0(R,\sig)$ by its module $\KL_{-1}^\s(R,\sig)$. For instance, for $R$ a Dedekind domain whose fraction field is a global field, we obtain that $\KL^\s_0(R[\ZZ],\bar{\id})$ is isomorphic to the split square-zero extension $\W^\s(R) \oplus \Pic(R)/2$, see Corollaries~\refthree{corollary:compute-L-groups-global-field} and~\refthree{corollary:compute-L-groups-global-Dedekind-ring-characteristic-2}.
\end{example}


\begin{remark}
Let $(R,\sig)$ be a regular noetherian ring with involution for which $2 \neq 0$ holds. Then $\KGW^\s(R,\sig) = \GW^\s(R,\sig)$ is not complex orientable: indeed, in $\pi_1$ of any complex orientable ring spectrum the image of the Hopf element $\eta \in \pi_1\SS$ vanishes, whereas 
the above assumptions %
imply that this image is non-zero in $\K_1(R)$, hence already in $\GW_1^\s(R,\sig)$.
In particular, Lemma~\reffour{lemma:sq0} implies that the $\Ek{0}$-equivalence $\GW^\s(R[\ZZ],\bar{\sig}) \simeq \GW^\s(R,\sig) \oplus \Sig\GW^\s(R,\sig)$ does not extend to one of $\Eone$-algebras. Nevertheless, in this situation,
$\GW_0^\s(R[\ZZ],\bar{\sig})$ is isomorphic to the square-zero extension $\GW_0^\s(R,\sig) \oplus \GW_{-1}^\s(R,\sig)$. Indeed, to see this, it suffices to argue that the augmentation ideal $\GW_{-1}^\s(R,\sig)$ has trivial multiplication. But the map $\GW_{-1}^\s(R,\sig) \to \L_{-1}^\s(R,\sig)$ is an isomorphism and is in addition compatible with its multiplication since $\GW \to \L$ is lax symmetric monoidal. Hence, the claim follows from Remark~\reffour{remark:L-group-sq0}. 
\end{remark}















\subsection{Multiplicative structures on normal L-spectra}
\label{subsection:normal-L}%

The purpose of the present section is to record that for any $\Einf$-Poincar\'e $\infty$-category $(\C,\QF)$ the associated map $\L(\C,\QF^\qdr) \rightarrow \L(\C,\QF)$ exhibits the source as a Smith ideal in the target, and consequently that its cofibre, the normal $\L$-spectrum inherits a multiplicative structure. To set this up, first observe that for any Poincar\'e $\infty$-category $(\C,\QF)$ we have a cube 
\[
\begin{tikzcd}
[column sep=0.5em, row sep=0.7em]
  & \L(\C,\QF^\qdr)\ar[rr]\ar[dd] && \KL(\C,\QF^\qdr)\ar[dd]\\
\GW(\C,\QF^\qdr)\ar[ru]\ar[rr, crossing over]\ar[dd] && \KGW(\C,\QF^\qdr)\ar[ur]\\
  & \L(\C,\QF)\ar[rr] && \KL(\C,\QF)\\
\GW(\C,\QF)\ar[rr] \ar[ru] && \KGW(\C,\QF)\ar[uu,<-, crossing over] \ar[ru]\\
\end{tikzcd}
\]
whose front, back and side faces are exact by Theorem \reffour{theorem:Karoubi-pullback-squares} and Corollary~\reftwo{corollary:tate-square-L}.
In particular, we find a square of equivalences
\[
\begin{tikzcd}
\cof(\GW(\C,\QF^\qdr) \rightarrow \GW(\C,\QF)) \ar[r] \ar[d] & \cof(\L(\C,\QF^\qdr) \rightarrow \L(\C,\QF)) \ar[d] \\
\cof(\KGW(\C,\QF^\qdr) \rightarrow \KGW(\C,\QF)) \ar[r] & \cof(\KL(\C,\QF^\qdr) \rightarrow \KL(\C,\QF))
\end{tikzcd}
\]
natural in $(\C,\QF) \in \Catp$.

\begin{definition}
We denote by $\NL \colon \Catp \rightarrow \Sp$, \emph{normal $\L$-theory}, the functor given by any of the four formulae in the square above.
\end{definition}

This notion was first introduced by Ranicki.
We will often use the abbreviation $\NL^r(R,M)$ for $\NL(\Dperf(R),\QF^r_M)$, where $M$ is an invertible module with involution over $R$.

Normal $\L$-spectra are closely connected to real topological cyclic homology for Poincar\'e $\infty$-categories and will be more thoroughly analysed in the forthcoming \cite{Harpaz-Nikolaus-Shah}, one of the main results of which equates them with the geometric fixed points of $\mathrm{TCR}(\C,\QF)$ under suitable assumptions.
Here, we will equip $\NL$ with a lax symmetric monoidal structure under $\KL \colon \Catp \rightarrow \Sp$ (and thus also under $\GW,\KGW$ and $\L$). 
To this end, we first recall that the pushout-tensor product 
\[
(e \rightarrow e') \otimes (f \rightarrow f') = [e\otimes f' \amalg_{e \otimes f} e' \otimes f \rightarrow e' \otimes f']
\]
on the arrow category of a symmetric monoidal $\infty$-category $\E$ with finite colimits is nothing but Day convolution 
with respect to the symmetric monoidal structure on $[1]$ given by the minimum (or equivalently, by the product). We will denote the corresponding symmetric monoidal $\infty$-category by 
\[
\Ar(\E)^\Sm := \Fun([1],\E)^\Day.
\]
Its unit is the unique arrow $\emptyset \to \one$ from the initial object to the unit of $\E$, and its algebras %
are precisely the Smith ideal inclusions.
There is also the pointwise symmetric monoidal structure on arrows, and it will be denoted by $\Ar(\E)^{\mathrm{pw}}$; it is similarly given by Day convolution with respect to the maximum (or equivalently, the coproduct). 


\begin{proposition}
\label{proposition:qdr-lax-monoidal}%
The functor 
\[
\Catp \longrightarrow \Ar(\Catp)^{\Sm}, \quad (\C,\QF) \longmapsto [(\C,\QF^\qdr) \rightarrow (\C,\QF)]
\]
admits a canonical lax symmetric monoidal structure.
\end{proposition}  
\begin{proof}
For $(\C,\QF) \in \Catp$, the natural map 
\[
(\C,\QF^{\qdr}) \to (\C,\QF) \otimes (\Spaf,\QF^\qdr),
\]
adjoint to the equivalence on underlying stable $\infty$-categories with duality on both sides under the adjunction of Proposition~\refone{proposition:adjunction-for-poincare-categories},
is an equivalence by the explicit formula for the tensor product in $\Catp$.
In then follows that the object
\[
[(\Spaf,\QF^\qdr) \to (\Spaf,\QF^\uni)] \in \Ar(\Catp)
\]
is idempotent as an $\Ek{0}$-algebra in $\Ar(\Catp)^{\Sm}$. Hence, it essentially uniquely refines to an idempotent $\Einf$-algebra and tensoring with it determines a smashing localisation $L \colon \Ar(\Catp)^{\Sm} \to \Ar(\Catp)^{\Sm}$ which is canonically a lax symmetric monoidal functor. Furthermore, the left Kan extension along the symmetric monoidal functor $\{1\} \to [1]$ induces a symmetric monoidal functor $\Catp \to \Ar(\Catp)^{\Sm}$. The composite of these two functors is then given by
\[
\Catp \to \Ar(\Catp)^{\Sm} \xrightarrow{L} \Ar(\Catp)^{\Sm}, \quad (\C,\QF) \mapsto (0\to (\C,\QF)) \mapsto (\C,\QF^\qdr) \to (\C,\QF),
\]
and so the proof is complete.
\end{proof}

We thus find that for any lax symmetric monoidal $\F \colon \Catp \rightarrow \Sp$ the assignment
\[
(\C,\QF) \longmapsto [\F(\C,\QF^\qdr) \rightarrow \F(\C,\QF)]
\]
upgrades to a lax symmetric monoidal functor $\Catp \rightarrow \Ar(\Sp)^{\Sm}$, in fact naturally in $\F$. In particular, for any $(\C,\QF) \in \Alg_{\Einf}(\Catp)$ we find that $\F(\C,\QF^\qdr)$ is a Smith ideal in the $\Einf$-ring $\F(\C,\QF)$. Finally, we note that taking the cofibre of an arrow induces a symmetric monoidal equivalence
\[
\cof \colon \Ar(\Sp)^{\Sm} \to \Ar(\Sp)^{\mathrm{pw}},
\]
see e.g.\ the argument for \cite{almostpaper}*{Proposition 2.3}.
We obtain:
\begin{corollary}
\label{corollary:normallmult}%
The tautological transformation $\KL \Rightarrow \NL$ of functors $\Catp \rightarrow \Sp$ admits a canonical lax symmetric monoidal structure extending that on the source constructed in Proposition \reffour{corollary:kl-monoidal}. Moreover, $\NL$ is Karoubi-localising and bordism invariant.
\end{corollary}
\begin{proof}
The first part follows immediately from Proposition~\reffour{proposition:qdr-lax-monoidal}. Then $\NL$ is a module over $\KL$ and hence bordism invariant. Finally, the functor $(\C,\QF) \mapsto (\C,\QF^\qdr)$ preserves Poincar\'e-Karoubi squares, for instance since it is given by the functor $- \otimes (\Spaf,\QF^\qdr)$ as noted above, so we may appeal to Proposition~\reffour{prop:closure-properties-of-Rs}. It follows that $\NL$ is Karoubi-localising as it is a cofibre of two Karoubi-localising functors.
\end{proof}


Ranicki already observed (see, e.g., \cite{RanickiTSO}) that normal $\L$-spectra of commutative rings (based on the symmetric Poincar\'e structures) admit homotopy associative and commutative multiplications. An upgrade regarding $\Einf$-structures on normal L-spectra was envisioned by Lurie \cite{Lurie-L-theory}*{Lecture 37}. 




\begin{example}
\label{example:F2}%
Perhaps the simplest non-trivial example of a normal $\L$-spectrum is $\NL^{\sym}(\FF_2)$, which is a 2-periodic ring spectrum all of whose homotopy groups are $\FF_2$. 
Since $\FF_2$ has cohomological dimension $0$, Corollaries~\refthree{corollary:genuine-is-quadratic} and \refthree{corollary:genuine-is-symmetric} together implies that
\[
\NL^{\geq m}(\FF_2) \simeq \tau_{\geq 2m-1}\NL^{\sym}(\FF_2) .
\]
In particular, for $m=0$ we get that the $(-1)$-connective cover map
\[
\NL^{\gs}(\FF_2) = \tau_{\geq -1}\NL^{\sym}(\FF_2) \to \NL^{\sym}(\FF_2)
\]
is an $\Einf$-ring spectrum map. %
The ring structure on $\NL^{\sym}_*(\FF_2)$ is thus
\[
\NL^{\sym}_*(\FF_2) = \FF_2[\beta,\beta^{-1},\eps]/\eps^2 = \L^{\sym}_*(\FF_2)[\eps]/\eps^2 \quad\text{and}\quad \NL^{\gs}_*(\FF_2) = \FF_2[\beta,\eps]/\eps^2 = \L^{\gs}_*(\FF_2)[\eps]/\eps^2
\]
with $|\beta| = 2$ and $|\eps| =-1$.
In particular, $\NL^{\sym}(\FF_2)$ has the same homotopy ring as $\L^{\sym}(\FF_2)^{S^1}$, though it is not a-priori clear if and why there would be any map relating the two.
This point is elucidated in forthcoming work \cite{Harpaz-Nikolaus-Shah}, where it is shown that there is always a canonical map $\L(\C,\QF)^{S^1} \to \NL(\C,\QF)$. In the case of $\FF_2$, this map induces equivalences of $\Einf$-rings
\[
\NL^{\sym}(\FF_2) \simeq \L^{\sym}(\FF_2)^{S^1} \quad\text{and}\quad \NL^{\gs}(\FF_2) \simeq \L^{\gs}(\FF_2)^{S^1} .
\]
Using the results of Heuts--Land \cite{formality}*{Theorem A \& C}, it follows that $\NL^\gs(\FF_2)$ and $\NL^\s(\FF_2)$ are split square zero extensions as augmented $\Eone$-algebras, but not
as augmented $\Ek{2}$-$\L^\gs(\FF_2)$-algebras. 
\end{example}



\begin{example}
\label{example:normal-L-of-Z}%
In the case of the integers, the normal symmetric $\L$-groups are given by
\[
\NL_i^{\sym}(\mathbb Z) \cong \begin{cases} \mathbb Z/8 & i \equiv 0 \ (4) \\ \mathbb Z/2 & i \equiv 1 \ (2) \\ 0 & i \equiv 2 \ (4)\end{cases}
\]
and the ring structure is $\NL^{\sym}_\ast(\mathbb Z) = \mathbb Z/8[t^{\pm 1},e,f]/(2e,2f,e^2,f^2,4-ef)$ where $|t|=4$, $|e| =1$ and $|f| = -1$, originally determined by Ranicki in unpublished work, see \cite{HLN}*{Appendix A} for a proof. 
Let us also recall from \cite{HLN}*{Corollary 6.2} that the map $\NL^\gs(\ZZ) \to \NL^\s(\ZZ)$ exhibits the source as the $(-1)$-connective cover of the target, just as is the case for $\FF_2$ recalled in Example~\reffour{example:F2}.
In particular, $\tau_{\geq-1}\NL^\s(\ZZ)$ admits an $\Einf$-ring structure, a phenomenon also originally observed in \cite{HLN}.

More generally, from Corollary~\refthree{corollary:genuine-is-quadratic} we find that $\NL^{\gs}(R,M)$ is $(-1)$-connective for every ordinary (in fact, connective $\Eone$-) ring $R$ and invertible module with involution $M$, and by Corollary~\refthree{corollary:genuine-is-symmetric} the resulting map
\[
\NL^{\gs}(R,M) \to \tau_{\geq -1}\NL^{\sym}(R,M),
\]
is an isomorphism in degrees $\geq d-1$ if $R$ is coherent of finite global dimension $d$. In a recent paper \cite{Schlichting-genuine}, Schlichting proved that if $R$ is a commutative Noetherian domain of finite Krull dimension either containing a field or smooth over a Dedekind domain and $M=R$ with the trivial $\Ct$-action, then the above map is in fact an equivalence.
In particular, for such $R$ %
we have that $\tau_{\geq -1}\NL^{\sym}(R)$ admits an $\Einf$-structure, yielding a non-trivial constraint on the $\Einf$-structure of $\NL^{\sym}(R)$. 
\end{example}


\begin{corollary}
\label{corollary:normalL8tors}%
Let $R$ be an $\Eone$-$\mathbb Z$-algebra
and $M$ an invertible module with involution over $R$. Then the normal $L$-spectra $\NL^{\geq m}(R,M)$ are naturally $\NL^{\gs}(\ZZ)$-modules for $m \in \ZZ \cup \{\infty,-\infty\}$. In particular, for every $i \in \ZZ$ and every $m,n \in \mathbb Z \cup \{\infty,- \infty\}$ with $m\geq n$, both kernel and cokernel of the natural maps
\[
\GW_i^{\geq m}(R,M) \longrightarrow \GW_i^{\geq n}(R,M)
\]
are $8$-torsion, and the same holds for any of the functors $\L,\KL,\KGW$ in place of $\GW$.
\end{corollary}
\begin{proof}
By Example~\reffour{example:ideal} and Corollary~\reffour{corollary:normallmult}, the normal L-spectra $\NL^{\geq m}(R,M)$ are module spectra over $\NL^{\gs}(\ZZ)$, and for $m\geq n$, the map $\NL^{\geq m}(R,M) \to \NL^{\geq n}(R,M)$ is one of $\NL^{\gs}(\ZZ)$-modules. Using then the equivalences
\begin{align*}
	\cof[\GW^{\geq m}(R,M) \to \GW^{\geq n}(R,M)] & \simeq \cof[\KGW^{\geq m}(R,M) \to \KGW^{\geq n}(R,M)] \\
		& \simeq \cof[\KL^{\geq m}(R,M) \to \KL^{\geq n}(R,M)] \\
		& \simeq \cof[\L^{\geq m}(R,M) \to \L^{\geq n}(R,M)] \\
		& \simeq \cof[\NL^{\geq m}(R,M) \to \NL^{\geq n}(R,M)]
\end{align*}
we find that all of these objects carry the structure of a $\NL^\gs(\ZZ)$-module. Hence the claim follows from the computation $\NL^\gs_0(\ZZ) \cong \ZZ/8$, see Example~\reffour{example:normal-L-of-Z}.
\end{proof}

Finally, we record the following invariance property of normal L-theory of $\ZZ$-algebras.
\begin{example}
Let $(A,\sig)$ be a $\ZZ$-algebra spectrum with involution in which 2 is invertible and $m\in \ZZ \cup \{\pm \infty \}$. Then $\NL^{\geq m}(A,\sig)$ vanishes since $\L^\qdr(A,\sig) \to \L^{\geq m}(A,\sig)$ is an equivalence. As a consequence, if $(R,\sig)$ is a discrete ring and $m\in \ZZ \cup \{\pm \infty \}$, the maps 
\[
\NL^{\geq m}(R,\sig) \to \NL^{\geq m}(R_{(2)},\sig) \to \NL^{\geq m}(R_2^\cwedge,\sig)
\]
are equivalences, that is, genuine normal L-spectra of $\ZZ$-algebras are insensitive to localising or (derived) completing at 2; this was implicitly already proven and used in \paperthree. Indeed, since $\NL$ is Karoubi-localising, this follows from the Poincar\'e-Karoubi square
\[
\begin{tikzcd}
	(\Dperf(R),\QF^{\geq m}_\sig) \ar[r] \ar[d] & (\Dperf(R_2^\cwedge),\QF^{\geq m}_\sig) \ar[d] \\
	(\Dperf(R[\tfrac{1}{2}],\QF^{\geq m}_\sig) \ar[r] & (\Dperf((R_2^\cwedge)[\tfrac{1}{2}],\QF^{\geq m}_{\sig})
\end{tikzcd}
\]
established in Proposition~\refthree{proposition:LC} and the previous claim. Note that the apparent assumption about the boundedness of the 2-torsion in $R$ in loc.\ cit.\ was only made to ensure that the derived completion $R^\cwedge_2$, for which the above Poincar\'e-Karoubi square always holds, is a discrete ring.
\end{example}












\subsection{The automorphism rings of Grothendieck-Witt- and L-theory}
\label{subsection:automorphism-GW-L}%

Finally, we use Proposition~\reftwo{proposition:pi0gwuniv} to determine the automorphisms of the Grothendieck-Witt and $\L$-theory functors. Yoneda's lemma, the universal property of the Grothendieck-Witt spectrum, Corollary~\reffour{corollary:GW-universal-multiplicative} and Proposition~\refone{proposition:corepresentability-of-poinc} provide equivalences
\[
\Nat(\GW,\GW) \simeq \Nat(\Poinc,\Omega^\infty\GW) \simeq \Nat(\Hom_{\Catp}((\Spa^\perf,\QF^\uni), -), \GWspace) \simeq \GWspace^\uni(\SS),
\]
with the $\Eone$-structure induced by composition on the left hand side which corresponds to the $\Eone$-structure on $\GWspace^\uni(\SS)$ induced by the underlying monoidal structure of $\GWspace$ and the symmetric monoidal structure of $(\Spa^\perf,\QF^\uni)$ as the unit of $\Catp$. Similarly,
$
\Nat(\L,\L) \simeq \Lspace^\uni(\SS),
$
as $\Eone$-spaces and bordification induces the map
\[
\Nat(\GW,\GW) \longrightarrow \Nat(\L,\L),
\]
which identifies with
\[
\GWspace^\uni(\SS) \xrightarrow{\bord} \Lspace^\uni(\SS).
\]
Applying Proposition~\reftwo{proposition:pi0gwuniv} which identifies this map on $\pi_0$ we find:
\begin{corollary}
\label{corollary:automorphisms-GW-L}%
The above identifications provide isomorphisms
\[
\pi_0\Nat(\GW,\GW) \cong \ZZ[e,h]/I \quad \text{and} \quad \pi_0\Nat(\L,\L) \cong \ZZ[e]/(e^2-8e)
\]
with $I = (e^2-8e,he-8h,h^2-2h)$ as in Proposition~\reftwo{proposition:pi0gwuniv}. In particular, we have 
\[
\pi_0\Aut(\GW) = \{\pm \id, \pm (\id-\hyp \circ \fgt)\} \cong (\Ct)^2 \quad \text{and} \quad \pi_0\Aut(\L) = \{\pm \id\} \cong \Ct.
\]
\end{corollary}


Finally, we note that this computation passes to the Karoubi versions of $\GW$ and $\L$:
\begin{corollary}
The canonical maps $\GW \to \KGW$ and $\L \rightarrow \KL$ induce equivalences
\[
\Nat(\GW,\GW) \simeq \Nat(\KGW,\KGW) \quad \text{and} \quad \Nat(\L,\L) \simeq \Nat(\KL,\KL)
\]
\end{corollary}
\begin{proof}
By the universal properties of the functors involved and Proposition~\refone{proposition:corepresentability-of-poinc}, the statement is equivalent to the maps 
\[
\Omega^\infty\GW^\uni(\SS) \longrightarrow \Om^\infty\KGW^\uni(\SS) \quad \text{and} \quad \Omega^\infty\L^\uni(\SS) \longrightarrow \Omega^\infty \KL^\uni(\SS)
\]
being equivalences, which is a consequence of Corollary~\reffour{corollary:regular-noetherian}.
\end{proof}
